\documentclass[8pt]{beamer}
\usetheme{Madrid}
\usecolortheme{default}
\usepackage{amsmath, amssymb, graphicx, array, multirow, booktabs, xcolor, tikz}
\usetikzlibrary{shapes,arrows,positioning,fit,backgrounds}

% Compact font settings
\setbeamerfont{title}{size=\Large}
\setbeamerfont{subtitle}{size=\small}
\setbeamerfont{author}{size=\scriptsize}
\setbeamerfont{date}{size=\scriptsize}
\setbeamerfont{frametitle}{size=\large}
\setbeamerfont{framesubtitle}{size=\normalsize}
\setbeamerfont{enumerate item}{size=\footnotesize}
\setbeamerfont{itemize item}{size=\footnotesize}
\setbeamerfont{block title}{size=\footnotesize}

% Compact spacing
\setlength{\itemsep}{0.08ex}
\setlength{\parskip}{0.08ex}
\setlength{\leftmargini}{0.3cm}
\setbeamersize{text margin left=3mm, text margin right=3mm}
\addtobeamertemplate{frame begin}{\vspace{-0.4cm}}{}

% Colors
\definecolor{myblue}{RGB}{0,0,255}
\definecolor{myred}{RGB}{255,0,0}
\definecolor{mygreen}{RGB}{0,128,0}
\definecolor{myorange}{RGB}{255,165,0}
\definecolor{mypurple}{RGB}{128,0,128}
\definecolor{mygray}{RGB}{128,128,128}
\definecolor{mygold}{RGB}{255,215,0}

\title{Module 4: Responsible \& Ethical AI}
\subtitle{100 Tutorial Topics: Ethical Frameworks, Principles, Bias, Privacy, Governance, Regulation}
\author{GARP RAI Exam Preparation}
\date{}

\begin{document}
	
	\begin{frame}
		\titlepage
	\end{frame}
	
	% ============================================================
	% SECTION 1: INTRODUCTION TO RESPONSIBLE AI
	% ============================================================
	
	\begin{frame}{T1: Responsible AI Definition}
		\fontsize{6.5}{7.5}\selectfont
		\textbf{TOPIC: What is Responsible AI?}
		
		\vspace{0.1cm}
		\textbf{DEFINITION:}
		\begin{itemize}
			\item Responsible AI refers to the ethical and responsible development, deployment, and use of artificial intelligence systems
			\item Aims to ensure that AI benefits society holistically while mitigating potential risks
			\item Involves aligning AI systems in a variety of ways to ensure they operate in a fair, transparent, and accountable manner
		\end{itemize}
		
		\vspace{0.1cm}
		\textbf{ALIGNMENT AREAS:}
		\begin{itemize}
			\item \textcolor{myblue}{Ethical Standards:} Moral principles and values
			\item \textcolor{myblue}{Industry Standards:} Best practices and guidelines
			\item \textcolor{myblue}{Regulation:} Legal requirements and compliance
			\item \textcolor{myblue}{Legislation:} Statutory laws
			\item \textcolor{myblue}{Global Principles:} International norms and agreements
		\end{itemize}
		
		\vspace{0.1cm}
		\textbf{KEY GOAL:}
		\begin{itemize}
			\item Benefit society holistically
			\item Mitigate potential risks
			\item Ensure fair, transparent, accountable operation
		\end{itemize}
		
		\textcolor{myred}{\textbf{EXAM TRAP:}} Responsible AI = \textcolor{myblue}{ethical development and deployment} of AI systems!
	\end{frame}
	
	\begin{frame}{T2: Purpose of Ethical Frameworks}
		\fontsize{6.5}{7.5}\selectfont
		\textbf{TOPIC: Why Ethical Frameworks Matter in AI}
		
		\vspace{0.1cm}
		\textbf{PURPOSE:}
		\begin{itemize}
			\item Help us decide how to design AI
			\item Help us deploy AI in a way that minimizes risks, harms, and wrongs
			\item Shape AI by helping define principles to abide by
		\end{itemize}
		
		\vspace{0.1cm}
		\textbf{KEY FUNCTIONS:}
		\begin{enumerate}
			\item \textcolor{myblue}{Define Principles:} What to abide by
			\item \textcolor{myblue}{Identify Biases:} Unacceptable biases
			\item \textcolor{myblue}{Assess Transparency:} What kind is desirable
			\item \textcolor{myblue}{Identify Stakeholders:} And their interests
			\item \textcolor{myblue}{Promote Accountability:} Who is responsible
			\item \textcolor{myblue}{Justify Decisions:} About design and implementation
		\end{enumerate}
		
		\vspace{0.1cm}
		\textbf{KEY BENEFIT:}
		\begin{itemize}
			\item Different ethical frameworks help make justifiable decisions
			\item Particularly important in pluralistic societies
			\item People hold diverse values and beliefs
		\end{itemize}
		
		\textcolor{myred}{\textbf{EXAM TRAP:}} Ethical frameworks \textcolor{myblue}{guide decision-making} to minimize risks and harms!
	\end{frame}
	
	\begin{frame}{T3: Practical Ethics Definition}
		\fontsize{6.5}{7.5}\selectfont
		\textbf{TOPIC: What is Practical Ethics?}
		
		\vspace{0.1cm}
		\textbf{DEFINITION:}
		\begin{itemize}
			\item The application of ethical theory and the development of good practices
			\item Goal: Solve real-world moral dilemmas
			\item Provide concrete guidance for moral decision making and problem solving
		\end{itemize}
		
		\vspace{0.1cm}
		\textbf{KEY ASPECTS:}
		\begin{enumerate}
			\item \textcolor{myblue}{Applied Focus:}
			\begin{itemize}
				\item Addresses actual ethical dilemmas
				\item Domains: medicine, law, politics, business, AI
				\item Offers recommendations, not just theories
			\end{itemize}
			\item \textcolor{myblue}{Multidisciplinary:}
			\begin{itemize}
				\item Draws on moral philosophy
				\item Social sciences (psychology, sociology)
				\item Domain expertise, law, computer science
			\end{itemize}
			\item \textcolor{myblue}{Context-Sensitive:}
			\begin{itemize}
				\item Situational nuances matter
				\item No one-size-fits-all approaches
			\end{itemize}
			\item \textcolor{myblue}{Pluralistic:}
			\begin{itemize}
				\item Blends different ethical frameworks
				\item Aims for justifiable decisions
			\end{itemize}
		\end{enumerate}
		
		\textcolor{myred}{\textbf{EXAM TRAP:}} Practical ethics = \textcolor{myblue}{applying theory to real-world dilemmas}!
	\end{frame}
	
	\begin{frame}{T4: Business Benefits of Ethics}
		\fontsize{6.5}{7.5}\selectfont
		\textbf{TOPIC: Why Firms Should Consider Practical Ethics}
		
		\vspace{0.1cm}
		\textbf{KEY BENEFITS:}
		\begin{enumerate}
			\item \textcolor{myblue}{Building Trust and Reputation:}
			\begin{itemize}
				\item Customers, employees, partners trust ethical companies
				\item Improves brand image
				\item Creates competitive advantage
			\end{itemize}
			\item \textcolor{myblue}{Avoiding Scandals:}
			\begin{itemize}
				\item Unethical behavior damages reputation and bottom line
				\item Ethics helps prevent major scandals
			\end{itemize}
			\item \textcolor{myblue}{Attracting and Retaining Talent:}
			\begin{itemize}
				\item Workers increasingly care about ethical practices
				\item Helps recruit and retain talent
			\end{itemize}
			\item \textcolor{myblue}{Strengthening Culture:}
			\begin{itemize}
				\item Nurtures teamwork, accountability, integrity
				\item Boosts morale and productivity
			\end{itemize}
			\item \textcolor{myblue}{Supporting Risk Management:}
			\begin{itemize}
				\item Reduces risk of safety failures, data breaches, non-compliance
			\end{itemize}
		\end{enumerate}
		
		\textcolor{myred}{\textbf{EXAM TRAP:}} Ethics makes \textcolor{myblue}{good business sense} with multiple benefits!
	\end{frame}
	
	\begin{frame}{T5: Ethics and Law Relationship}
		\fontsize{6.5}{7.5}\selectfont
		\textbf{TOPIC: How Ethics and Law Relate}
		
		\vspace{0.1cm}
		\textbf{ETHICS AS COMPLEMENT TO LAW:}
		\begin{itemize}
			\item Ethics helps ground, inform, and shape law
			\item Societies regulate behavior based on moral acceptability
			\item Ethics helps distinguish between just and unjust laws
		\end{itemize}
		
		\vspace{0.1cm}
		\textbf{LAWS ARE NARROW IN SCOPE:}
		\begin{itemize}
			\item Establish minimal behavioral requirements
			\item Enable orderly interactions within basic fairness
			\item Provide a framework for society
		\end{itemize}
		
		\vspace{0.1cm}
		\textbf{ETHICS GOES BEYOND LAW:}
		\begin{itemize}
			\item Identifies moral issues
			\item Reflects on the kind of society we want
			\item Makes recommendations for a good life
			\item More ambitious than law
		\end{itemize}
		
		\vspace{0.1cm}
		\textbf{EXAMPLE:}
		\begin{itemize}
			\item It's not illegal to abandon friends, but it's immoral
			\item Laws provide minimums; ethics provides ideals
		\end{itemize}
		
		\textcolor{myred}{\textbf{EXAM TRAP:}} Ethics \textcolor{myblue}{complements and informs} law - goes beyond minimum requirements!
	\end{frame}
	
	% ============================================================
	% SECTION 2: ETHICAL FRAMEWORKS - CONSEQUENTIALISM
	% ============================================================
	
	\begin{frame}{T6: Consequentialism Definition}
		\fontsize{6.5}{7.5}\selectfont
		\textbf{TOPIC: Understanding Consequentialism}
		
		\vspace{0.1cm}
		\textbf{DEFINITION:}
		\begin{itemize}
			\item An ethical theory that judges the morality of an action based on the consequences of that action
			\item Focuses on outcomes or results to determine whether an action is right or wrong
		\end{itemize}
		
		\vspace{0.1cm}
		\textbf{KEY PRINCIPLES:}
		\begin{enumerate}
			\item \textcolor{myblue}{Outcome Focus:}
			\begin{itemize}
				\item What results from the action?
				\item Does it produce good or bad outcomes?
				\item Consequences determine moral value
			\end{itemize}
			\item \textcolor{myblue}{Forward-Looking:}
			\begin{itemize}
				\item Focus on future results
				\item Past actions judged by their effects
				\item Decisions based on expected outcomes
			\end{itemize}
			\item \textcolor{myblue}{Circumstantially Relative:}
			\begin{itemize}
				\item What's right depends on circumstances
				\item No absolute rules
				\item Context matters
			\end{itemize}
		\end{enumerate}
		
		\textcolor{myred}{\textbf{EXAM TRAP:}} Consequentialism = \textcolor{myblue}{judging actions by their consequences}!
	\end{frame}
	
	\begin{frame}{T7: Utilitarianism}
		\fontsize{6.5}{7.5}\selectfont
		\textbf{TOPIC: Utilitarianism - The Most Common Form of Consequentialism}
		
		\vspace{0.1cm}
		\textbf{CORE IDEA:}
		\begin{itemize}
			\item Maximize overall utility
			\item Utility often defined as happiness, pleasure, or satisfaction of desires
			\item The morally correct action produces the greatest net utility for all affected
		\end{itemize}
		
		\vspace{0.1cm}
		\textbf{KEY THINKERS:}
		\begin{enumerate}
			\item \textcolor{myblue}{Jeremy Bentham:}
			\begin{itemize}
				\item Founder of utilitarianism
				\item Hedonic calculus
				\item Focus on pleasure and pain
			\end{itemize}
			\item \textcolor{myblue}{John Stuart Mill:}
			\begin{itemize}
				\item Refined utilitarianism
				\item Quality of pleasures
				\item Higher and lower pleasures
			\end{itemize}
		\end{enumerate}
		
		\vspace{0.1cm}
		\textbf{APPLICATION TO AI:}
		\begin{itemize}
			\item An AI system is morally good if it increases overall pleasure/happiness
			\item Actions are not intrinsically moral
			\item Derive moral value solely from results
		\end{itemize}
		
		\textcolor{myred}{\textbf{EXAM TRAP:}} Utilitarianism = \textcolor{myblue}{maximize overall utility/happiness}!
	\end{frame}
	
	\begin{frame}{T8: Consequentialism Advantages}
		\fontsize{6.5}{7.5}\selectfont
		\textbf{TOPIC: Advantages of Consequentialism}
		
		\vspace{0.1cm}
		\textbf{KEY ADVANTAGES:}
		\begin{enumerate}
			\item \textcolor{myblue}{Single Quantifiable Metric:}
			\begin{itemize}
				\item Provides one measure of moral value
				\item Utility can be calculated
				\item Clear decision-making framework
			\end{itemize}
			\item \textcolor{myblue}{Impartiality:}
			\begin{itemize}
				\item Everyone's utility counts equally
				\item No special treatment
				\item Objective calculation
			\end{itemize}
			\item \textcolor{myblue}{Scientific Amenability:}
			\begin{itemize}
				\item Utility can be measured
				\item Evidence-based decisions
				\item Quantifiable outcomes
			\end{itemize}
			\item \textcolor{myblue}{Forward-Looking:}
			\begin{itemize}
				\item Focuses on future outcomes
				\item Practical decision-making
				\item Real-world consequences
			\end{itemize}
		\end{enumerate}
		
		\textcolor{myred}{\textbf{EXAM TRAP:}} Consequentialism provides \textcolor{myblue}{a single quantifiable metric} for moral decisions!
	\end{frame}
	
	\begin{frame}{T9: Consequentialism Challenges}
		\fontsize{6.5}{7.5}\selectfont
		\textbf{TOPIC: Criticisms of Consequentialism}
		
		\vspace{0.1cm}
		\textbf{KEY CHALLENGES:}
		\begin{enumerate}
			\item \textcolor{myblue}{Subjectivity of Utility:}
			\begin{itemize}
				\item Choice of metric is subjective
				\item Different definitions of "good"
				\item Measurement difficulties
			\end{itemize}
			\item \textcolor{myblue}{Unpredictable Consequences:}
			\begin{itemize}
				\item Can't predict all outcomes
				\item Where does the calculation stop?
				\item Immediate vs. long-term consequences
			\end{itemize}
			\item \textcolor{myblue}{Rights Violations:}
			\begin{itemize}
				\item Can justify violating individual rights
				\item "The ends justify the means"
				\item Classic example: Killing one to save five
			\end{itemize}
			\item \textcolor{myblue}{Doctor Example:}
			\begin{itemize}
				\item Killing a patient to use organs saves five
				\item Maximizes utility but seems immoral
				\item Challenge for consequentialism
			\end{itemize}
		\end{enumerate}
		
		\textcolor{myred}{\textbf{EXAM TRAP:}} Consequentialism can justify \textcolor{myblue}{violating rights for the greater good}!
	\end{frame}
	
	\begin{frame}{T10: Rule Consequentialism}
		\fontsize{6.5}{7.5}\selectfont
		\textbf{TOPIC: Rule Consequentialism - A Modification}
		
		\vspace{0.1cm}
		\textbf{CORE IDEA:}
		\begin{itemize}
			\item Judges acts by whether they adhere to utility-maximizing rules
			\item Not just by case-specific outcomes
			\item Addresses issues with act utilitarianism
		\end{itemize}
		
		\vspace{0.1cm}
		\textbf{HOW IT WORKS:}
		\begin{enumerate}
			\item \textcolor{myblue}{Follow General Rules:}
			\begin{itemize}
				\item Rules that tend to maximize utility
				\item Rules against murder, lying, etc.
				\item Upheld even when breaking them may increase utility
			\end{itemize}
			\item \textcolor{myblue}{Example - Doctor Case:}
			\begin{itemize}
				\item Doctor should not kill patient
				\item Would undermine trust in medical system
				\item Long-term utility is maximized by following rules
			\end{itemize}
			\item \textcolor{myblue}{AI Application:}
			\begin{itemize}
				\item Should not violate rights even if efficient
				\item Long-term trust is valuable
				\item Rule-based protection
			\end{itemize}
		\end{enumerate}
		
		\textcolor{myred}{\textbf{EXAM TRAP:}} Rule consequentialism = \textcolor{myblue}{follow utility-maximizing rules}!
	\end{frame}
	
	% ============================================================
	% SECTION 3: ETHICAL FRAMEWORKS - DEONTOLOGY
	% ============================================================
	
	\begin{frame}{T11: Deontology Definition}
		\fontsize{6.5}{7.5}\selectfont
		\textbf{TOPIC: Understanding Deontology}
		
		\vspace{0.1cm}
		\textbf{DEFINITION:}
		\begin{itemize}
			\item An ethical framework that judges the morality of actions based on adherence to ethical duties and rules
			\item Rather than focusing on consequences
			\item Focus on moral obligations regardless of outcomes
		\end{itemize}
		
		\vspace{0.1cm}
		\textbf{KEY PRINCIPLES:}
		\begin{enumerate}
			\item \textcolor{myblue}{Duty-Based:}
			\begin{itemize}
				\item Moral obligations are paramount
				\item We must do our duty
				\item Consequences don't override duties
			\end{itemize}
			\item \textcolor{myblue}{Rule-Focused:}
			\begin{itemize}
				\item Follow moral rules
				\item Universal principles
				\item Consistency across situations
			\end{itemize}
			\item \textcolor{myblue}{Respect for Persons:}
			\begin{itemize}
				\item People are ends in themselves
				\item Never treat as mere means
				\item Intrinsic value of individuals
			\end{itemize}
		\end{enumerate}
		
		\textcolor{myred}{\textbf{EXAM TRAP:}} Deontology = \textcolor{myblue}{judging actions by duties and rules}!
	\end{frame}
	
	\begin{frame}{T12: Immanuel Kant}
		\fontsize{6.5}{7.5}\selectfont
		\textbf{TOPIC: Kant's Contribution to Deontology}
		
		\vspace{0.1cm}
		\textbf{KEY THINKER:}
		\begin{itemize}
			\item Immanuel Kant is the most important proponent of deontology
			\item Emphasized moral principles, rationality, and duty
			\item Famous for not placing moral weight on consequences
			\item Believed lying is always wrong regardless of consequences
		\end{itemize}
		
		\vspace{0.1cm}
		\textbf{CORE CONCEPTS:}
		\begin{enumerate}
			\item \textcolor{myblue}{Categorical Imperative:}
			\begin{itemize}
				\item Act only according to maxims that could be universal law
				\item "It's wrong to lie" - applicable universally
				\item No contradictions in universal application
			\end{itemize}
			\item \textcolor{myblue}{Treating People as Ends:}
			\begin{itemize}
				\item Never treat people as mere means to an end
				\item People have intrinsic value
				\item Respect their autonomy and dignity
			\end{itemize}
			\item \textcolor{myblue}{Reason-Based Morality:}
			\begin{itemize}
				\item Morality grounded in reason
				\item Not in emotions or consequences
				\item Rational beings can determine moral duties
			\end{itemize}
		\end{enumerate}
		
		\textcolor{myred}{\textbf{EXAM TRAP:}} Kant = \textcolor{myblue}{duty-based ethics focused on universal principles}!
	\end{frame}
	
	\begin{frame}{T13: Categorical Imperative}
		\fontsize{6.5}{7.5}\selectfont
		\textbf{TOPIC: Understanding the Categorical Imperative}
		
		\vspace{0.1cm}
		\textbf{FIRST FORMULATION:}
		\begin{itemize}
			\item "Act only according to that maxim whereby you can at the same time will that it should become a universal law"
			\item Can the rule be applied universally without contradiction?
			\item Example: Could "It's okay to lie" be a universal law?
		\end{itemize}
		
		\vspace{0.1cm}
		\textbf{SECOND FORMULATION:}
		\begin{itemize}
			\item "Act in such a way that you treat humanity, whether in your own person or in the person of any other, always at the same time as an end and never merely as a means"
			\item Respect the dignity of all persons
			\item Don't use people as tools
		\end{itemize}
		
		\vspace{0.1cm}
		\textbf{APPLICATION TO AI:}
		\begin{itemize}
			\item AI systems should respect human dignity
			\item People should not be treated as data points
			\item Transparency and consent are essential
		\end{itemize}
		
		\textcolor{myred}{\textbf{EXAM TRAP:}} Categorical imperative = \textcolor{myblue}{act on rules that could be universal law}!
	\end{frame}
	
	\begin{frame}{T14: Modern Deontology}
		\fontsize{6.5}{7.5}\selectfont
		\textbf{TOPIC: Contemporary Deontological Ethics}
		
		\vspace{0.1cm}
		\textbf{EVOLUTION OF DEONTOLOGY:}
		\begin{itemize}
			\item Most contemporary deontologists care about consequences
			\item However, consequences are not the only moral consideration
			\item There are red lines that should not be crossed
		\end{itemize}
		
		\vspace{0.1cm}
		\textbf{KEY FEATURES:}
		\begin{enumerate}
			\item \textcolor{myblue}{Rights:}
			\begin{itemize}
				\item Individuals have fundamental rights
				\item Rights cannot be violated even for good outcomes
				\item Examples: Privacy, autonomy, non-discrimination
			\end{itemize}
			\item \textcolor{myblue}{Rules:}
			\begin{itemize}
				\item Established ethical rules matter
				\item Professional codes of conduct
				\item Organizational policies
			\end{itemize}
			\item \textcolor{myblue}{Principles:}
			\begin{itemize}
				\item Ethical principles provide guidance
				\item Beyond consequences
				\item Deontological in nature
			\end{itemize}
		\end{enumerate}
		
		\textcolor{myred}{\textbf{EXAM TRAP:}} Modern deontology \textcolor{myblue}{considers consequences but has red lines}!
	\end{frame}
	
	\begin{frame}{T15: Deontology Challenges}
		\fontsize{6.5}{7.5}\selectfont
		\textbf{TOPIC: Criticisms of Deontology}
		
		\vspace{0.1cm}
		\textbf{KEY CHALLENGES:}
		\begin{enumerate}
			\item \textcolor{myblue}{Rigidity:}
			\begin{itemize}
				\item Can lead to rigid rule-following
				\item Even when consequences are disastrous
				\item Example: Kant said lying is always wrong
			\end{itemize}
			\item \textcolor{myblue}{Conflicting Duties:}
			\begin{itemize}
				\item What if duties conflict?
				\item Example: Duty to tell truth vs. duty to protect
				\item No clear resolution framework
			\end{itemize}
			\item \textcolor{myblue}{Application Difficulty:}
			\begin{itemize}
				\item Hard to apply in complex situations
				\item What rules apply?
				\item How to resolve conflicts?
			\end{itemize}
		\end{enumerate}
		
		\vspace{0.1cm}
		\textbf{AI APPLICATION:}
		\begin{itemize}
			\item Can lead to overly rigid AI systems
			\item May not adapt to context
			\item Difficulty encoding all rules
		\end{itemize}
		
		\textcolor{myred}{\textbf{EXAM TRAP:}} Deontology can be \textcolor{myblue}{too rigid} even when consequences are disastrous!
	\end{frame}
	
	% ============================================================
	% SECTION 4: ETHICAL FRAMEWORKS - VIRTUE ETHICS
	% ============================================================
	
	\begin{frame}{T16: Virtue Ethics Definition}
		\fontsize{6.5}{7.5}\selectfont
		\textbf{TOPIC: Understanding Virtue Ethics}
		
		\vspace{0.1cm}
		\textbf{DEFINITION:}
		\begin{itemize}
			\item Emphasizes virtuous character traits and living a good life
			\item Rather than rules or consequences
			\item Has roots in ancient Greek philosophers like Aristotle
		\end{itemize}
		
		\vspace{0.1cm}
		\textbf{KEY PRINCIPLES:}
		\begin{enumerate}
			\item \textcolor{myblue}{Character Focus:}
			\begin{itemize}
				\item What kind of person should I be?
				\item Cultivating virtues
				\item Becoming a good person
			\end{itemize}
			\item \textcolor{myblue}{Virtues:}
			\begin{itemize}
				\item Wisdom, courage, humanity
				\item Justice, temperance, generosity
				\item Character traits to cultivate
			\end{itemize}
			\item \textcolor{myblue}{Practical Wisdom:}
			\begin{itemize}
				\item Exercising judgment in situations
				\item Moderating emotions and appetites
				\item Appropriate behavior in context
			\end{itemize}
		\end{enumerate}
		
		\textcolor{myred}{\textbf{EXAM TRAP:}} Virtue ethics = \textcolor{myblue}{focus on character and living a good life}!
	\end{frame}
	
	\begin{frame}{T17: Aristotle and Virtue Ethics}
		\fontsize{6.5}{7.5}\selectfont
		\textbf{TOPIC: Aristotle's Contribution to Virtue Ethics}
		
		\vspace{0.1cm}
		\textbf{HISTORICAL ROOTS:}
		\begin{itemize}
			\item Virtue ethics has roots in ancient Greek philosophers
			\item Aristotle taught that happiness comes from a life guided by virtues
			\item Virtues are nurtured through practice and habit
		\end{itemize}
		
		\vspace{0.1cm}
		\textbf{KEY CONCEPTS:}
		\begin{enumerate}
			\item \textcolor{myblue}{Eudaimonia:}
			\begin{itemize}
				\item Flourishing or living well
				\item Ultimate goal of human life
				\item Achieved through virtue
			\end{itemize}
			\item \textcolor{myblue}{Virtues as Habits:}
			\begin{itemize}
				\item Virtues are developed through practice
				\item Repeated actions create character
				\item Modeling virtuous exemplars
			\end{itemize}
			\item \textcolor{myblue}{Golden Mean:}
			\begin{itemize}
				\item Virtue lies between extremes
				\item Courage between cowardice and recklessness
				\item Balance in character
			\end{itemize}
		\end{enumerate}
		
		\textcolor{myred}{\textbf{EXAM TRAP:}} Aristotle = \textcolor{myblue}{virtues through practice for a good life}!
	\end{frame}
	
	\begin{frame}{T18: Key Virtues}
		\fontsize{6.5}{7.5}\selectfont
		\textbf{TOPIC: Examples of Virtues}
		
		\vspace{0.1cm}
		\textbf{CORE VIRTUES:}
		\begin{enumerate}
			\item \textcolor{myblue}{Wisdom:}
			\begin{itemize}
				\item Practical judgment
				\item Understanding what matters
				\item Sound decision-making
			\end{itemize}
			\item \textcolor{myblue}{Courage:}
			\begin{itemize}
				\item Facing challenges appropriately
				\item Moral courage
				\item Standing up for what's right
			\end{itemize}
			\item \textcolor{myblue}{Humanity:}
			\begin{itemize}
				\item Compassion and care
				\item Empathy for others
				\item Benevolence
			\end{itemize}
			\item \textcolor{myblue}{Justice:}
			\begin{itemize}
				\item Fairness and equity
				\item Giving proper due
				\item Social responsibility
			\end{itemize}
			\item \textcolor{myblue}{Temperance:}
			\begin{itemize}
				\item Self-control
				\item Moderation
				\item Balance
			\end{itemize}
			\item \textcolor{myblue}{Generosity:}
			\begin{itemize}
				\item Giving freely
				\item Sharing resources
				\item Beneficence
			\end{itemize}
		\end{enumerate}
		
		\textcolor{myred}{\textbf{EXAM TRAP:}} Virtues include \textcolor{myblue}{wisdom, courage, justice, temperance, and generosity}!
	\end{frame}
	
	\begin{frame}{T19: Virtue Ethics in Practice}
		\fontsize{6.5}{7.5}\selectfont
		\textbf{TOPIC: Applying Virtue Ethics}
		
		\vspace{0.1cm}
		\textbf{PRACTICAL APPROACH:}
		\begin{enumerate}
			\item \textcolor{myblue}{Ask the Virtue Question:}
			\begin{itemize}
				\item "What would a virtuous person do?"
				\item Consider role models and exemplars
				\item What would Jesus, Mohammed, Gandhi do?
			\end{itemize}
			\item \textcolor{myblue}{Cultivate Virtues:}
			\begin{itemize}
				\item Through practice and habit
				\item Learning from experience
				\item Modeling virtuous behavior
			\end{itemize}
			\item \textcolor{myblue}{Consider Character:}
			\begin{itemize}
				\item Not just discrete acts
				\item Life-long character development
				\item Holistic approach to morality
			\end{itemize}
		\end{enumerate}
		
		\vspace{0.1cm}
		\textbf{APPLICATION TO AI:}
		\begin{itemize}
			\item Build AI that learns from experience like children
			\item Develop virtuous organizational culture
			\item Consider what responsible companies would do
		\end{itemize}
		
		\textcolor{myred}{\textbf{EXAM TRAP:}} Virtue ethics asks \textcolor{myblue}{"What would a virtuous person do?"}!
	\end{frame}
	
	\begin{frame}{T20: Virtue Ethics Challenges}
		\fontsize{6.5}{7.5}\selectfont
		\textbf{TOPIC: Criticisms of Virtue Ethics}
		
		\vspace{0.1cm}
		\textbf{KEY CHALLENGES:}
		\begin{enumerate}
			\item \textcolor{myblue}{Lack of Clear Guidance:}
			\begin{itemize}
				\item Less concrete than rules-based approaches
				\item Hard to know what virtue requires
				\item Situational judgment needed
			\end{itemize}
			\item \textcolor{myblue}{Conflicting Virtues:}
			\begin{itemize}
				\item Different virtues can conflict
				\item How to weigh courage vs. prudence?
				\item No clear resolution method
			\end{itemize}
			\item \textcolor{myblue}{Cultural Relativism:}
			\begin{itemize}
				\item Virtues vary across cultures
				\item What's virtuous in one context may not be in another
				\item Need for universal standards
			\end{itemize}
		\end{enumerate}
		
		\vspace{0.1cm}
		\textbf{VIRTUE ETHICIST RESPONSE:}
		\begin{itemize}
			\item Practical wisdom helps navigate hard cases
			\item Other theories also have conflicts
			\item Consequences not always comparable
		\end{itemize}
		
		\textcolor{myred}{\textbf{EXAM TRAP:}} Virtue ethics criticized for \textcolor{myblue}{lack of clear guidance}!
	\end{frame}
	
	% ============================================================
	% SECTION 5: FRAMEWORK COMPLEMENTARITY AND MEDICAL ETHICS
	% ============================================================
	
	\begin{frame}{T21: Framework Complementarity}
		\fontsize{6.5}{7.5}\selectfont
		\textbf{TOPIC: How Ethical Frameworks Complement Each Other}
		
		\vspace{0.1cm}
		\textbf{COMBINED APPROACH:}
		\begin{itemize}
			\item Consequentialism, deontology, and virtue ethics are not mutually exclusive
			\item In practical ethics, they complement one another
			\item The best moral decision accords with all three
		\end{itemize}
		
		\vspace{0.1cm}
		\textbf{INTEGRATED DECISION-MAKING:}
		\begin{enumerate}
			\item \textcolor{myblue}{Consequentialist Analysis:}
			\begin{itemize}
				\item What are the likely outcomes?
				\item Does it maximize good consequences?
				\item Consider all affected parties
			\end{itemize}
			\item \textcolor{myblue}{Deontological Analysis:}
			\begin{itemize}
				\item Does it respect rights and duties?
				\item Are there red lines being crossed?
				\item Is it in accordance with principles?
			\end{itemize}
			\item \textcolor{myblue}{Virtue Ethics Analysis:}
			\begin{itemize}
				\item What would a virtuous person do?
				\item Does it build character?
				\item Is it consistent with values?
			\end{itemize}
		\end{enumerate}
		
		\textcolor{myred}{\textbf{EXAM TRAP:}} Frameworks are \textcolor{myblue}{complementary, not mutually exclusive}!
	\end{frame}
	
	\begin{frame}{T22: Medical Ethics History}
		\fontsize{6.5}{7.5}\selectfont
		\textbf{TOPIC: Medical Ethics - A Historical Overview}
		
		\vspace{0.1cm}
		\textbf{ANCIENT FOUNDATIONS:}
		\begin{itemize}
			\item Hippocratic Oath (5th-3rd century BC)
			\item "Do no harm" to patients
			\item First known code of medical ethics
		\end{itemize}
		
		\vspace{0.1cm}
		\textbf{KEY DOCUMENTS:}
		\begin{enumerate}
			\item \textcolor{myblue}{Nuremberg Code (1947):}
			\begin{itemize}
				\item Response to Nazi medical experiments
				\item Informed consent requirement
				\item Research ethics foundation
			\end{itemize}
			\item \textcolor{myblue}{Declaration of Geneva (1948):}
			\begin{itemize}
				\item Modern Hippocratic Oath
				\item International medical ethics
			\end{itemize}
			\item \textcolor{myblue}{Declaration of Helsinki (1964):}
			\begin{itemize}
				\item Research ethics guidelines
				\item Updated regularly
				\item Global standards
			\end{itemize}
			\item \textcolor{myblue}{Belmont Report (1978):}
			\begin{itemize}
				\item Respect for persons, beneficence, justice
				\item Research ethics framework
				\item US foundational document
			\end{itemize}
		\end{enumerate}
		
		\textcolor{myred}{\textbf{EXAM TRAP:}} Medical ethics has \textcolor{myblue}{long history and established frameworks}!
	\end{frame}
	
	\begin{frame}{T23: Medical Ethics Development}
		\fontsize{6.5}{7.5}\selectfont
		\textbf{TOPIC: Medical Ethics Development in the 1970s}
		
		\vspace{0.1cm}
		\textbf{FACTORS DRIVING EVOLUTION:}
		\begin{enumerate}
			\item \textcolor{myblue}{Concentration of Medical Care:}
			\begin{itemize}
				\item Hospitals and depersonalized settings
				\item Loss of personal relationships
				\item Need for formal ethics
			\end{itemize}
			\item \textcolor{myblue}{Rising Costs and Government Role:}
			\begin{itemize}
				\item Increased government involvement
				\item Healthcare funding changes
				\item Resource allocation decisions
			\end{itemize}
			\item \textcolor{myblue}{Patient Rights Movements:}
			\begin{itemize}
				\item Civil rights outgrowth
				\item Demands for autonomy
				\item Informed consent
			\end{itemize}
			\item \textcolor{myblue}{Medical Scandals:}
			\begin{itemize}
				\item Tuskegee Syphilis Experiment
				\item Public outrage
				\item Call for reform
			\end{itemize}
			\item \textcolor{myblue}{Rapid Technological Advances:}
			\begin{itemize}
				\item Mechanical ventilators
				\item Organ transplantation
				\item New ethical dilemmas
			\end{itemize}
		\end{enumerate}
		
		\textcolor{myred}{\textbf{EXAM TRAP:}} Multiple factors drove \textcolor{myblue}{development of medical ethics}!
	\end{frame}
	
	\begin{frame}{T24: Lessons for AI Ethics}
		\fontsize{6.5}{7.5}\selectfont
		\textbf{TOPIC: What AI Ethics Can Learn from Medical Ethics}
		
		\vspace{0.1cm}
		\textbf{KEY LESSONS:}
		\begin{enumerate}
			\item \textcolor{myblue}{Structured Frameworks:}
			\begin{itemize}
				\item Medical ethics has established principles
				\item AI needs similar foundations
				\item Clear ethical guidance
			\end{itemize}
			\item \textcolor{myblue}{Professional Ethics:}
			\begin{itemize}
				\item Doctors have ethical training
				\item AI professionals need ethics education
				\item Professional responsibility
			\end{itemize}
			\item \textcolor{myblue}{Ethics Committees:}
			\begin{itemize}
				\item Medical ethics committees
				\item AI needs similar oversight
				\item Multi-disciplinary review
			\end{itemize}
			\item \textcolor{myblue}{Codes of Ethics:}
			\begin{itemize}
				\item Medical codes are well-established
				\item AI codes are evolving
				\item Need for standards
			\end{itemize}
		\end{enumerate}
		
		\textcolor{myred}{\textbf{EXAM TRAP:}} Medical ethics provides \textcolor{myblue}{model for AI ethics development}!
	\end{frame}
	
	% ============================================================
	% SECTION 6: PRINCIPLES OF AI ETHICS
	% ============================================================
	
	\begin{frame}{T25: Five Core Principles Overview}
		\fontsize{6.5}{7.5}\selectfont
		\textbf{TOPIC: Five Core Principles of AI Ethics}
		
		\vspace{0.1cm}
		\textbf{THE FIVE PRINCIPLES:}
		\begin{enumerate}
			\item \textcolor{myblue}{Nonmaleficence:}
			\begin{itemize}
				\item Obligation to avoid harming others
				\item "Do no harm"
				\item Prohibition of unjustified risks
			\end{itemize}
			\item \textcolor{myblue}{Beneficence:}
			\begin{itemize}
				\item Obligation to act for the benefit of others
				\item "Do good"
				\item Active help and support
			\end{itemize}
			\item \textcolor{myblue}{Justice:}
			\begin{itemize}
				\item Fairness, equality, impartiality
				\item Equitable distribution
				\item Non-discrimination
			\end{itemize}
			\item \textcolor{myblue}{Autonomy:}
			\begin{itemize}
				\item Respect for informed, un-coerced decisions
				\item Individual control
				\item Non-manipulation
			\end{itemize}
			\item \textcolor{myblue}{Explainability:}
			\begin{itemize}
				\item Transparent and understandable decisions
				\item Accountability
				\item Trust-building
			\end{itemize}
		\end{enumerate}
		
		\textcolor{myred}{\textbf{EXAM TRAP:}} Five principles = \textcolor{myblue}{nonmaleficence, beneficence, justice, autonomy, explainability}!
	\end{frame}
	
	\begin{frame}{T26: Nonmaleficence}
		\fontsize{6.5}{7.5}\selectfont
		\textbf{TOPIC: The Principle of Nonmaleficence}
		
		\vspace{0.1cm}
		\textbf{CORE MEANING:}
		\begin{itemize}
			\item Obligation to avoid harming others
			\item Prohibition on imposing unjustified risks
			\item "First, do no harm" (primum non nocere)
		\end{itemize}
		
		\vspace{0.1cm}
		\textbf{KEY ASPECTS:}
		\begin{enumerate}
			\item \textcolor{myblue}{Avoid Actual Harm:}
			\begin{itemize}
				\item Don't directly harm others
				\item Example: Algorithms that discriminate
				\item Actions that cause injury
			\end{itemize}
			\item \textcolor{myblue}{Avoid Unjustified Risk:}
			\begin{itemize}
				\item Don't impose unnecessary risks
				\item Example: Untested algorithms
				\item Risks that outweigh benefits
			\end{itemize}
			\item \textcolor{myblue}{Proportionality:}
			\begin{itemize}
				\item Some risk is permissible
				\item If benefits justify risk
				\item If alternatives unavailable
			\end{itemize}
		\end{enumerate}
		
		\textcolor{myred}{\textbf{EXAM TRAP:}} Nonmaleficence = \textcolor{myblue}{"do no harm"} - avoid harming others!
	\end{frame}
	
	\begin{frame}{T27: Risk Justification}
		\fontsize{6.5}{7.5}\selectfont
		\textbf{TOPIC: When is Risk Acceptable?}
		
		\vspace{0.1cm}
		\textbf{PRINCIPLES OF RISK ACCEPTANCE:}
		\begin{enumerate}
			\item \textcolor{myblue}{Benefit Justification:}
			\begin{itemize}
				\item Risk is acceptable if benefits justify it
				\item Example: Clinical trials
				\item Risk is part of medical progress
			\end{itemize}
			\item \textcolor{myblue}{No Alternatives:}
			\begin{itemize}
				\item Risk acceptable if no alternatives
				\item Example: Emergency procedures
				\item Only option available
			\end{itemize}
			\item \textcolor{myblue}{Who Bears the Risk:}
			\begin{itemize}
				\item More acceptable if those at risk benefit
				\item Example: Clinical trial participants
				\item Unacceptable if others bear risk
			\end{itemize}
		\end{enumerate}
		
		\vspace{0.1cm}
		\textbf{AI APPLICATION:}
		\begin{itemize}
			\item Imposing algorithmic risk on people who don't benefit
			\item Unacceptable to experiment on populations
			\item Especially if they don't gain
		\end{itemize}
		
		\textcolor{myred}{\textbf{EXAM TRAP:}} Risk acceptable if \textcolor{myblue}{benefits justify it and risk-bearers benefit}!
	\end{frame}
	
	\begin{frame}{T28: Beneficence}
		\fontsize{6.5}{7.5}\selectfont
		\textbf{TOPIC: The Principle of Beneficence}
		
		\vspace{0.1cm}
		\textbf{CORE MEANING:}
		\begin{itemize}
			\item Moral obligation to act for the benefit of others
			\item Requires taking positive steps to help
			\item Moves beyond simply refraining from harm
		\end{itemize}
		
		\vspace{0.1cm}
		\textbf{KEY ASPECTS:}
		\begin{enumerate}
			\item \textcolor{myblue}{Active Help:}
			\begin{itemize}
				\item Donating to charity
				\item Volunteering
				\item Providing resources and assistance
			\end{itemize}
			\item \textcolor{myblue}{Beyond Nonmaleficence:}
			\begin{itemize}
				\item Not just "do no harm"
				\item "Do good"
				\item Active beneficence
			\end{itemize}
			\item \textcolor{myblue}{Deontological Duty:}
			\begin{itemize}
				\item Duty-based obligation
				\item Not just about maximizing utility
				\item Required beyond consequences
			\end{itemize}
		\end{enumerate}
		
		\textcolor{myred}{\textbf{EXAM TRAP:}} Beneficence = \textcolor{myblue}{"do good"} - actively help others!
	\end{frame}
	
	\begin{frame}{T29: Beneficence Limitations}
		\fontsize{6.5}{7.5}\selectfont
		\textbf{TOPIC: Limits to the Duty of Beneficence}
		
		\vspace{0.1cm}
		\textbf{REASONABLE CONSTRAINTS:}
		\begin{enumerate}
			\item \textcolor{myblue}{Scarce Resources:}
			\begin{itemize}
				\item Cannot help everyone
				\item Limited resources
				\item Prioritization needed
			\end{itemize}
			\item \textcolor{myblue}{Competing Obligations:}
			\begin{itemize}
				\item Multiple duties to balance
				\item Professional responsibilities
				\item Personal obligations
			\end{itemize}
			\item \textcolor{myblue}{Reasonableness:}
			\begin{itemize}
				\item There's only so much ethics can demand
				\item Demandingness constraint
				\item Proportionality matters
			\end{itemize}
			\item \textcolor{myblue}{Recipient Autonomy:}
			\begin{itemize}
				\item Unwanted "benefits" may be rejected
				\item Right to decide what's best
				\item Respect individual choices
			\end{itemize}
		\end{enumerate}
		
		\textcolor{myred}{\textbf{EXAM TRAP:}} Beneficence has \textcolor{myblue}{reasonable limits} - cannot help everyone!
	\end{frame}
	
	\begin{frame}{T30: Justice}
		\fontsize{6.5}{7.5}\selectfont
		\textbf{TOPIC: The Principle of Justice}
		
		\vspace{0.1cm}
		\textbf{CORE MEANING:}
		\begin{itemize}
			\item Fairness, equality, impartiality, proportionality
			\item Giving each person their proper due
			\item Upholding duties toward fairness and equality
		\end{itemize}
		
		\vspace{0.1cm}
		\textbf{KEY APPLICATIONS:}
		\begin{enumerate}
			\item \textcolor{myblue}{Fair Treatment:}
			\begin{itemize}
				\item No discrimination based on protected characteristics
				\item Equal opportunities
				\item Impartial algorithms
			\end{itemize}
			\item \textcolor{myblue}{Equitable Distribution:}
			\begin{itemize}
				\item Benefits and burdens fairly allocated
				\item Access to AI benefits
				\item Protection from AI harms
			\end{itemize}
			\item \textcolor{myblue}{Accountability:}
			\begin{itemize}
				\item Those responsible should be accountable
				\item Just processes
				\item Fair consequences
			\end{itemize}
		\end{enumerate}
		
		\textcolor{myred}{\textbf{EXAM TRAP:}} Justice = \textcolor{myblue}{fairness, equality, impartiality, and proportionality}!
	\end{frame}
	
	\begin{frame}{T31: Types of Justice}
		\fontsize{6.5}{7.5}\selectfont
		\textbf{TOPIC: Four Types of Justice}
		
		\vspace{0.1cm}
		\textbf{FOUR TYPES:}
		\begin{enumerate}
			\item \textcolor{myblue}{Procedural Justice:}
			\begin{itemize}
				\item Fair processes and impartiality
				\item Right structures in place
				\item Rule of law
			\end{itemize}
			\item \textcolor{myblue}{Distributive Justice:}
			\begin{itemize}
				\item Equitable allocation of benefits and burdens
				\item Fair distribution of resources
				\item Access to opportunities
			\end{itemize}
			\item \textcolor{myblue}{Restorative Justice:}
			\begin{itemize}
				\item Repairing harms
				\item Reconciling victims and offenders
				\item Making amends
			\end{itemize}
			\item \textcolor{myblue}{Interactional Justice:}
			\begin{itemize}
				\item Respect and fairness between individuals
				\item Interpersonal treatment
				\item Dignity and respect
			\end{itemize}
		\end{enumerate}
		
		\textcolor{myred}{\textbf{EXAM TRAP:}} Four types of justice = \textcolor{myblue}{procedural, distributive, restorative, interactional}!
	\end{frame}
	
	\begin{frame}{T32: Autonomy}
		\fontsize{6.5}{7.5}\selectfont
		\textbf{TOPIC: The Principle of Autonomy}
		
		\vspace{0.1cm}
		\textbf{CORE MEANING:}
		\begin{itemize}
			\item Capacity to make informed, un-coerced decisions
			\item Respecting others' abilities to determine their own course
			\item Non-manipulation
		\end{itemize}
		
		\vspace{0.1cm}
		\textbf{KEY ASPECTS:}
		\begin{enumerate}
			\item \textcolor{myblue}{Informed Decisions:}
			\begin{itemize}
				\item People need adequate information
				\item Transparency about options
				\item Understanding consequences
			\end{itemize}
			\item \textcolor{myblue}{Un-coerced Decisions:}
			\begin{itemize}
				\item No coercion or manipulation
				\item Free choice
				\item Voluntary consent
			\end{itemize}
			\item \textcolor{myblue}{Capacity:}
			\begin{itemize}
				\item Must have capacity to decide
				\item Exceptions: children, unconscious patients
				\item Surrogate decision makers when needed
			\end{itemize}
		\end{enumerate}
		
		\textcolor{myred}{\textbf{EXAM TRAP:}} Autonomy = \textcolor{myblue}{informed, un-coerced decisions}!
	\end{frame}
	
	\begin{frame}{T33: Autonomy in Healthcare}
		\fontsize{6.5}{7.5}\selectfont
		\textbf{TOPIC: Autonomy in Medical Context}
		
		\vspace{0.1cm}
		\textbf{INFORMED CONSENT:}
		\begin{itemize}
			\item Patients have right to voluntary informed consent
			\item Right to refuse treatment
			\item Doctor's duty to inform
			\item Patient's right to decide
		\end{itemize}
		
		\vspace{0.1cm}
		\textbf{WHAT UNDERMINES AUTONOMY:}
		\begin{enumerate}
			\item \textcolor{myblue}{Coercion:}
			\begin{itemize}
				\item Forcing decisions
				\item Threats or pressure
				\item Removing choice
			\end{itemize}
			\item \textcolor{myblue}{Deception:}
			\begin{itemize}
				\item Lying or misleading
				\item Withholding information
				\item Manipulating understanding
			\end{itemize}
			\item \textcolor{myblue}{Undue Influence:}
			\begin{itemize}
				\item Excessive persuasion
				\item Exploitation
				\item Imbalance of power
			\end{itemize}
		\end{enumerate}
		
		\textcolor{myred}{\textbf{EXAM TRAP:}} Autonomy requires \textcolor{myblue}{informed consent and no coercion}!
	\end{frame}
	
	\begin{frame}{T34: Autonomy in AI}
		\fontsize{6.5}{7.5}\selectfont
		\textbf{TOPIC: Autonomy in the Context of AI}
		
		\vspace{0.1cm}
		\textbf{ETHICAL AI AND AUTONOMY:}
		\begin{itemize}
			\item Ethical AI respects people's autonomy
			\item No coercive or manipulative tactics
			\item Technology should help people further their own goals
			\item Not serve third-party interests
		\end{itemize}
		
		\vspace{0.1cm}
		\textbf{RISKS TO AUTONOMY:}
		\begin{enumerate}
			\item \textcolor{myblue}{Manipulation:}
			\begin{itemize}
				\item Recommendation algorithms
				\item Targeted content
				\item Behavioral influence
			\end{itemize}
			\item \textcolor{myblue}{Coercion:}
			\begin{itemize}
				\item Forced compliance
				\item No alternative
				\item System lock-in
			\end{itemize}
			\item \textcolor{myblue}{Lack of Transparency:}
			\begin{itemize}
				\item Users don't understand decisions
				\item Cannot make informed choices
				\item Autonomy undermined
			\end{itemize}
		\end{enumerate}
		
		\textcolor{myred}{\textbf{EXAM TRAP:}} Ethical AI respects \textcolor{myblue}{autonomy by avoiding coercion and manipulation}!
	\end{frame}
	
	\begin{frame}{T35: Explainability}
		\fontsize{6.5}{7.5}\selectfont
		\textbf{TOPIC: The Principle of Explainability}
		
		\vspace{0.1cm}
		\textbf{CORE MEANING:}
		\begin{itemize}
			\item AI systems should provide transparent explanations
			\item Understandable reasons for actions or decisions
			\item Also called "explicability"
		\end{itemize}
		
		\vspace{0.1cm}
		\textbf{WHY EXPLAINABILITY MATTERS:}
		\begin{enumerate}
			\item \textcolor{myblue}{Accountability:}
			\begin{itemize}
				\item Hold designers and companies accountable
				\item Identify errors and biases
				\item Rectify unfair practices
			\end{itemize}
			\item \textcolor{myblue}{Trust:}
			\begin{itemize}
				\item Trust is fundamental for adoption
				\item Users trust understandable systems
				\item Transparency builds confidence
			\end{itemize}
			\item \textcolor{myblue}{Ethical Compliance:}
			\begin{itemize}
				\item Ensure adherence to ethical guidelines
				\item Respect individual rights
				\item Regulatory compliance
			\end{itemize}
		\end{enumerate}
		
		\textcolor{myred}{\textbf{EXAM TRAP:}} Explainability = \textcolor{myblue}{transparent and understandable explanations}!
	\end{frame}
	
	\begin{frame}{T36: Types of Explainability}
		\fontsize{6.5}{7.5}\selectfont
		\textbf{TOPIC: Levels and Types of Explainability}
		
		\vspace{0.1cm}
		\textbf{FOUR TYPES:}
		\begin{enumerate}
			\item \textcolor{myblue}{Local Explainability:}
			\begin{itemize}
				\item Explains decisions on a single instance
				\item "Why was THIS decision made?"
				\item Example: Why was THIS loan denied?
			\end{itemize}
			\item \textcolor{myblue}{Global Explainability:}
			\begin{itemize}
				\item Explains overall model behavior
				\item "How does the model work generally?"
				\item Comprehensive understanding
			\end{itemize}
			\item \textcolor{myblue}{Model-Specific Explainability:}
			\begin{itemize}
				\item Techniques tailored to architecture
				\item Example: Decision tree rules
				\item Neural networks require different methods
			\end{itemize}
			\item \textcolor{myblue}{Model-Agnostic Explainability:}
			\begin{itemize}
				\item Works with any model
				\item Versatile and general
				\item LIME, SHAP are examples
			\end{itemize}
		\end{enumerate}
		
		\textcolor{myred}{\textbf{EXAM TRAP:}} Explainability has \textcolor{myblue}{different levels and types} for different needs!
	\end{frame}
	
	\begin{frame}{T37: Counterfactual Explanations}
		\fontsize{6.5}{7.5}\selectfont
		\textbf{TOPIC: Counterfactual Explanations}
		
		\vspace{0.1cm}
		\textbf{WHAT THEY ARE:}
		\begin{itemize}
			\item Hypothetical scenarios contrasting with actual decisions
			\item Show conditions for different outcomes
			\item Provide actionable insights
		\end{itemize}
		
		\vspace{0.1cm}
		\textbf{EXAMPLE - LOAN DENIAL:}
		\begin{itemize}
			\item Actual: Loan denied
			\item Counterfactual: "You would have been approved if:
			\begin{itemize}
				\item You had \$10,000 more in the bank
				\item OR you earned \$1,000 more per month
			\end{itemize}
		\end{itemize}
		
		\vspace{0.1cm}
		\textbf{BENEFITS:}
		\begin{enumerate}
			\item \textcolor{myblue}{Transparency:}
			\begin{itemize}
				\item Users understand decision factors
				\item See what influences outcomes
			\end{itemize}
			\item \textcolor{myblue}{Actionable Insights:}
			\begin{itemize}
				\item Know what to change
				\item Path to different outcome
				\item Empowerment
			\end{itemize}
			\item \textcolor{myblue}{Fairness:}
			\begin{itemize}
				\item Detect biases
				\item Understand disparate treatment
				\item Identify improvement areas
			\end{itemize}
		\end{enumerate}
		
		\textcolor{myred}{\textbf{EXAM TRAP:}} Counterfactual explanations show \textcolor{myblue}{what would lead to different outcomes}!
	\end{frame}
	
	% ============================================================
	% SECTION 7: BIAS, DISCRIMINATION, AND FAIRNESS
	% ============================================================
	
	\begin{frame}{T38: Bias in AI}
		\fontsize{6.5}{7.5}\selectfont
		\textbf{TOPIC: Understanding Bias in AI Systems}
		
		\vspace{0.1cm}
		\textbf{DEFINITION:}
		\begin{itemize}
			\item Systemic deviation in output or impact
			\item Compared to a desired norm or standard
			\item Algorithm is biased when it deviates from intended function
		\end{itemize}
		
		\vspace{0.1cm}
		\textbf{EXAMPLE - HIRING ALGORITHM:}
		\begin{itemize}
			\item Designed to identify best candidates
			\item But recommends based on race or sex
			\item Deviates from intended purpose
			\item This is problematic bias
		\end{itemize}
		
		\vspace{0.1cm}
		\textbf{KEY INSIGHT:}
		\begin{itemize}
			\item No algorithm is perfectly objective
			\item Every algorithm reflects embedded values
			\item Values shape what's considered valuable
		\end{itemize}
		
		\textcolor{myred}{\textbf{EXAM TRAP:}} AI bias = \textcolor{myblue}{systemic deviation from intended function}!
	\end{frame}
	
	\begin{frame}{T39: Justifiable vs Problematic Bias}
		\fontsize{6.5}{7.5}\selectfont
		\textbf{TOPIC: When is Bias Acceptable?}
		
		\vspace{0.1cm}
		\textbf{JUSTIFIABLE BIAS:}
		\begin{itemize}
			\item Part of well-designed AI
			\item Optimizes for relevant criteria
			\item Example: Loan eligibility prioritizing bank balance
			\item Relevant to loan repayment
		\end{itemize}
		
		\vspace{0.1cm}
		\textbf{PROBLEMATIC BIAS:}
		\begin{itemize}
			\item Results in unfairness
			\item Disadvantages for unjustifiable reasons
			\item Based on irrelevant characteristics
			\item Example: Race or sex discrimination
		\end{itemize}
		
		\vspace{0.1cm}
		\textbf{STATISTICALLY UNBIASED $\neq$ ETHICALLY ACCEPTABLE:}
		\begin{itemize}
			\item Even unbiased algorithms can cause harm
			\item Example: Charging exorbitant fees to everyone
			\item Statistically fair but ethically problematic
		\end{itemize}
		
		\textcolor{myred}{\textbf{EXAM TRAP:}} Not all biases are \textcolor{myblue}{problematic} - some are justifiable!
	\end{frame}
	
	\begin{frame}{T40: Sources of Bias Overview}
		\fontsize{6.5}{7.5}\selectfont
		\textbf{TOPIC: Four Primary Sources of Algorithmic Bias}
		
		\vspace{0.1cm}
		\textbf{1. Problem Specification:}
		\begin{itemize}
			\item Goals contain inherent problems
			\item Target variables miss real-world objectives
			\item Example: Healthcare expenditure proxy
		\end{itemize}
		
		\vspace{0.1cm}
		\textbf{2. Data:}
		\begin{itemize}
			\item Historical bias
			\item Sampling bias
			\item Selective labels problem
		\end{itemize}
		
		\vspace{0.1cm}
		\textbf{3. Modeling, Validation, and Design:}
		\begin{itemize}
			\item Optimization choices
			\item Model selection
			\item Subgroup consideration
		\end{itemize}
		
		\vspace{0.1cm}
		\textbf{4. Deployment:}
		\begin{itemize}
			\item Human deference
			\item Unintended effects
			\item Context changes
		\end{itemize}
		
		\textcolor{myred}{\textbf{EXAM TRAP:}} Bias can enter at \textcolor{myblue}{any stage} of AI development!
	\end{frame}
	
	\begin{frame}{T41: Problem Specification Bias}
		\fontsize{6.5}{7.5}\selectfont
		\textbf{TOPIC: Bias in Problem Specification}
		
		\vspace{0.1cm}
		\textbf{HOW IT OCCURS:}
		\begin{enumerate}
			\item \textcolor{myblue}{Complex Goal Operationalization:}
			\begin{itemize}
				\item Hard to capture real-world objectives
				\item Selected target variables may be poor proxies
				\item Example: Bank eliminating all risk
			\end{itemize}
			\item \textcolor{myblue}{Unintended Consequences:}
			\begin{itemize}
				\item Bank lends only to affluent
				\item Excludes others
				\item Fairness issue
			\end{itemize}
			\item \textcolor{myblue}{Proxy Problems:}
			\begin{itemize}
				\item Healthcare expenditure proxy
				\item Favored richest patients, not sickest
				\item Systematically biased
			\end{itemize}
		\end{enumerate}
		
		\vspace{0.1cm}
		\textbf{KEY LESSON:}
		\begin{itemize}
			\item Careful problem definition is essential
			\item Proxies must be thoughtfully chosen
			\item Goals should reflect true objectives
		\end{itemize}
		
		\textcolor{myred}{\textbf{EXAM TRAP:}} Problem specification can introduce \textcolor{myblue}{bias through poor operationalization}!
	\end{frame}
	
	\begin{frame}{T42: Data-Related Biases}
		\fontsize{6.5}{7.5}\selectfont
		\textbf{TOPIC: Biases in Data}
		
		\vspace{0.1cm}
		\textbf{HISTORICAL BIAS:}
		\begin{itemize}
			\item Data reflects past discrimination
			\item Example: Women excluded from banking
			\item Algorithm favors men
			\item Perpetuates historical patterns
		\end{itemize}
		
		\vspace{0.1cm}
		\textbf{SELECTIVE LABELS PROBLEM:}
		\begin{itemize}
			\item Missing counterfactual outcomes
			\item No data on rejected applicants
			\item Can't know if better candidates existed
			\item Perpetuates existing patterns
		\end{itemize}
		
		\vspace{0.1cm}
		\textbf{SAMPLING BIAS:}
		\begin{itemize}
			\item Non-random data sample
			\item Trends don't generalize
			\item Example: Data from young men only
			\item May not apply to women or elderly
		\end{itemize}
		
		\textcolor{myred}{\textbf{EXAM TRAP:}} Data biases include \textcolor{myblue}{historical, selective labels, and sampling bias}!
	\end{frame}
	
	\begin{frame}{T43: Modeling and Design Bias}
		\fontsize{6.5}{7.5}\selectfont
		\textbf{TOPIC: Bias in Modeling, Validation, and Design}
		
		\vspace{0.1cm}
		\textbf{HOW BIAS EMERGES:}
		\begin{enumerate}
			\item \textcolor{myblue}{Optimization Functions:}
			\begin{itemize}
				\item Choice of objective
				\item What is being optimized?
				\item May favor certain outcomes
			\end{itemize}
			\item \textcolor{myblue}{Regression Models:}
			\begin{itemize}
				\item Different models have different biases
				\item Choice matters
				\item Subgroup effects
			\end{itemize}
			\item \textcolor{myblue}{Subgroup Consideration:}
			\begin{itemize}
				\item How are subgroups handled?
				\item Are differences recognized?
				\item Fairness across groups
			\end{itemize}
			\item \textcolor{myblue}{Information Presentation:}
			\begin{itemize}
				\item How is data presented?
				\item Visual biases
				\item Interface effects
			\end{itemize}
		\end{enumerate}
		
		\textcolor{myred}{\textbf{EXAM TRAP:}} Modeling choices can \textcolor{myblue}{introduce bias} even with good data!
	\end{frame}
	
	\begin{frame}{T44: Deployment Bias}
		\fontsize{6.5}{7.5}\selectfont
		\textbf{TOPIC: Bias in Deployment}
		
		\vspace{0.1cm}
		\textbf{HOW IT OCCURS:}
		\begin{enumerate}
			\item \textcolor{myblue}{Human Deference:}
			\begin{itemize}
				\item People defer to algorithm suggestions
				\item Even when limitations known
				\item Convenience factor
			\end{itemize}
			\item \textcolor{myblue}{Reasons for Deference:}
			\begin{itemize}
				\item Time-saving
				\item Perceived objectivity
				\item Self-doubt
				\item Responsibility shielding
			\end{itemize}
			\item \textcolor{myblue}{Unintended Effects:}
			\begin{itemize}
				\item Individuals relinquish responsibility
				\item Blame shifting
				\item Behavior changes
			\end{itemize}
		\end{enumerate}
		
		\vspace{0.1cm}
		\textbf{KEY LESSON:}
		\begin{itemize}
			\item Oversight is essential
			\item Human judgment remains important
			\item Don't blindly trust automation
		\end{itemize}
		
		\textcolor{myred}{\textbf{EXAM TRAP:}} Deployment can introduce \textcolor{myblue}{bias through human deference} to algorithms!
	\end{frame}
	
	\begin{frame}{T45: Automation Deference}
		\fontsize{6.5}{7.5}\selectfont
		\textbf{TOPIC: Why People Defer to Automation}
		
		\vspace{0.1cm}
		\textbf{CONVENIENCE:}
		\begin{itemize}
			\item Time-saving
			\item Effort reduction
			\item Easy to follow
		\end{itemize}
		
		\vspace{0.1cm}
		\textbf{PERCEIVED OBJECTIVITY:}
		\begin{itemize}
			\item Algorithms appear more objective
			\item People know their own fallibilities
			\item Creates self-doubt
		\end{itemize}
		
		\vspace{0.1cm}
		\textbf{RESPONSIBILITY SHIELDING:}
		\begin{itemize}
			\item Deferring shares responsibility
			\item Going against algorithm exposes people
			\item Blame avoidance
		\end{itemize}
		
		\vspace{0.1cm}
		\textbf{CONSEQUENCES:}
		\begin{itemize}
			\item Unintended effects
			\item Reduced oversight
			\item Bias perpetuation
			\item Accountability gaps
		\end{itemize}
		
		\textcolor{myred}{\textbf{EXAM TRAP:}} People defer to algorithms for \textcolor{myblue}{convenience, objectivity perception, and responsibility shielding}!
	\end{frame}
	
	\begin{frame}{T46: Discrimination in Law}
		\fontsize{6.5}{7.5}\selectfont
		\textbf{TOPIC: When Bias Becomes Discrimination}
		
		\vspace{0.1cm}
		\textbf{LEGAL PERSPECTIVE:}
		\begin{itemize}
			\item Depends on jurisdiction
			\item When it disfavors based on protected characteristics
			\item Race, sex, ethnicity, age
			\item Other classifications protected by law
		\end{itemize}
		
		\vspace{0.1cm}
		\textbf{PROTECTED CHARACTERISTICS:}
		\begin{enumerate}
			\item Race
			\item Sex/Gender
			\item Ethnicity
			\item Age
			\item Religion
			\item Disability
			\item Sexual orientation
		\end{enumerate}
		
		\vspace{0.1cm}
		\textbf{KEY POINT:}
		\begin{itemize}
			\item Such disadvantages typically violate legal protections
			\item Proactive measures needed
			\item Continuous monitoring
		\end{itemize}
		
		\textcolor{myred}{\textbf{EXAM TRAP:}} Discrimination = \textcolor{myblue}{disfavoring based on protected characteristics} in violation of law!
	\end{frame}
	
	\begin{frame}{T47: Fairness in AI}
		\fontsize{6.5}{7.5}\selectfont
		\textbf{TOPIC: Understanding Fairness in AI}
		
		\vspace{0.1cm}
		\textbf{DEFINITION:}
		\begin{itemize}
			\item Absence of bias or preference
			\item Based on irrelevant characteristics
			\item Algorithm is fair when it doesn't exhibit problematic biases
		\end{itemize}
		
		\vspace{0.1cm}
		\textbf{TWO TYPES:}
		\begin{enumerate}
			\item \textcolor{myblue}{Group Fairness:}
			\begin{itemize}
				\item Statistical criteria across groups
				\item Demographics mirror population
				\item Example: 51\% female = 51\% loan recipients
			\end{itemize}
			\item \textcolor{myblue}{Individual Fairness:}
			\begin{itemize}
				\item Treating similar individuals similarly
				\item Challenge: Defining "similar"
				\item Focus on individual treatment
			\end{itemize}
		\end{enumerate}
		
		\textcolor{myred}{\textbf{EXAM TRAP:}} Fairness = \textcolor{myblue}{absence of problematic biases}!
	\end{frame}
	
	\begin{frame}{T48: Fairness Impossibility}
		\fontsize{6.5}{7.5}\selectfont
		\textbf{TOPIC: Why Fairness Can't Be Automated}
		
		\vspace{0.1cm}
		\textbf{MATHEMATICAL IMPOSSIBILITY:}
		\begin{itemize}
			\item When base rates are unequal
			\item Cannot satisfy all fairness measures simultaneously
			\item Demographic parity, predictive rate parity, equal opportunity
			\item Trade-offs are unavoidable
		\end{itemize}
		
		\vspace{0.1cm}
		\textbf{EXAMPLE - CRIMINAL RISK ASSESSMENT:}
		\begin{itemize}
			\item Men statistically more likely to commit certain crimes
			\item Algorithm assigns higher risk to man
			\item Man argues unfair treatment
			\item Woman argues unfair if assigned equal score
		\end{itemize}
		
		\vspace{0.1cm}
		\textbf{KEY POINT:}
		\begin{itemize}
			\item Fairness is a moral judgment
			\item Not purely mathematical
			\item Compromises and trade-offs needed
		\end{itemize}
		
		\textcolor{myred}{\textbf{EXAM TRAP:}} Fairness cannot be automated when \textcolor{myblue}{base rates are unequal}!
	\end{frame}
	
	\begin{frame}{T49: Procedural Fairness}
		\fontsize{6.5}{7.5}\selectfont
		\textbf{TOPIC: Procedural Fairness in AI}
		
		\vspace{0.1cm}
		\textbf{DEFINITION:}
		\begin{itemize}
			\item Fair processes and impartial procedures
			\item Not just fair outcomes
			\item Right structures in place
		\end{itemize}
		
		\vspace{0.1cm}
		\textbf{JUSTICE SYSTEM ANALOGY:}
		\begin{itemize}
			\item Procedural fairness: rule of law
			\item Outcome fairness: appropriate punishments
			\item Mistakes can be justifiable with fair processes
			\item Ways to right wrongs
		\end{itemize}
		
		\vspace{0.1cm}
		\textbf{AI APPLICATION:}
		\begin{itemize}
			\item Ethics committees
			\item Established procedures
			\item Appeal processes
			\item Accountability structures
		\end{itemize}
		
		\textcolor{myred}{\textbf{EXAM TRAP:}} Procedural fairness = \textcolor{myblue}{fair processes and impartial procedures}!
	\end{frame}
	
	\begin{frame}{T50: Avoiding Problematic Biases}
		\fontsize{6.5}{7.5}\selectfont
		\textbf{TOPIC: Strategies to Avoid Problematic Biases}
		
		\vspace{0.1cm}
		\textbf{KEY APPROACHES:}
		\begin{enumerate}
			\item \textcolor{myblue}{Technological Solutions:}
			\begin{itemize}
				\item Fairness toolkits (Aequitas, AIF360)
				\item Internal auditing systems
				\item Bias reports from providers
				\item But technology alone is insufficient
			\end{itemize}
			\item \textcolor{myblue}{Data Quality:}
			\begin{itemize}
				\item Diverse, updated, accurate data
				\item Representative samples
				\item Free from past discriminatory tendencies
			\end{itemize}
			\item \textcolor{myblue}{Auditing:}
			\begin{itemize}
				\item Algorithmic auditing is best approach
				\item Private companies offer services
				\item Independent assessment
			\end{itemize}
			\item \textcolor{myblue}{Ethical Committees:}
			\begin{itemize}
				\item Discuss biases
				\item Consider trade-offs
				\item Ensure procedural fairness
			\end{itemize}
		\end{enumerate}
		
		\textcolor{myred}{\textbf{EXAM TRAP:}} Bias mitigation requires \textcolor{myblue}{multiple approaches beyond technology}!
	\end{frame}
	
	% ============================================================
	% SECTION 8: PRIVACY AND CYBERSECURITY
	% ============================================================
	
	\begin{frame}{T51: Privacy in AI Context}
		\fontsize{6.5}{7.5}\selectfont
		\textbf{TOPIC: Understanding Privacy in AI}
		
		\vspace{0.1cm}
		\textbf{DEFINITION:}
		\begin{itemize}
			\item Privacy to extent others don't have access to personal information
			\item In reference to some person or institution
			\item In reference to some personal data point
		\end{itemize}
		
		\vspace{0.1cm}
		\textbf{WHY PRIVACY MATTERS:}
		\begin{enumerate}
			\item \textcolor{myblue}{Protection from Abuse:}
			\begin{itemize}
				\item More knowledge enables interference
				\item Power imbalances
				\item Control over individuals
			\end{itemize}
			\item \textcolor{myblue}{Autonomy:}
			\begin{itemize}
				\item Privacy supports self-determination
				\item Freedom from surveillance
				\item Personal space
			\end{itemize}
			\item \textcolor{myblue}{Dignity:}
			\begin{itemize}
				\item Respect for persons
				\item Not being reduced to data
				\item Human worth
			\end{itemize}
		\end{enumerate}
		
		\textcolor{myred}{\textbf{EXAM TRAP:}} Privacy = \textcolor{myblue}{lack of access to personal information}!
	\end{frame}
	
	\begin{frame}{T52: Data Security}
		\fontsize{6.5}{7.5}\selectfont
		\textbf{TOPIC: Understanding Data Security}
		
		\vspace{0.1cm}
		\textbf{DEFINITION:}
		\begin{itemize}
			\item Doing what is necessary to prevent unauthorized access
			\item Protecting data from attacks and theft
			\item Ensuring data integrity
		\end{itemize}
		
		\vspace{0.1cm}
		\textbf{KEY COMPONENTS:}
		\begin{enumerate}
			\item \textcolor{myblue}{Physical Security:}
			\begin{itemize}
				\item Hardware protection
				\item Storage device security
				\item Physical access controls
			\end{itemize}
			\item \textcolor{myblue}{Administrative Controls:}
			\begin{itemize}
				\item Access permissions
				\item Organizational policies
				\item Employee training
			\end{itemize}
			\item \textcolor{myblue}{Software Applications:}
			\begin{itemize}
				\item Encryption
				\item Authentication
				\item Monitoring systems
			\end{itemize}
		\end{enumerate}
		
		\textcolor{myred}{\textbf{EXAM TRAP:}} Data security = \textcolor{myblue}{preventing unauthorized access} to data!
	\end{frame}
	
	\begin{frame}{T53: Privacy as Ethical Issue}
		\fontsize{6.5}{7.5}\selectfont
		\textbf{TOPIC: Why Privacy is an Ethical Issue}
		
		\vspace{0.1cm}
		\textbf{WRONGS:}
		\begin{itemize}
			\item Violates moral and legal rights
			\item Treats people instrumentally
			\item Reduces people to data points
			\item Respect for autonomy
		\end{itemize}
		
		\vspace{0.1cm}
		\textbf{HARMS TO INDIVIDUALS:}
		\begin{itemize}
			\item Discrimination
			\item Blackmail
			\item Exposure and shaming
			\item Identity theft
		\end{itemize}
		
		\vspace{0.1cm}
		\textbf{HARMS TO INSTITUTIONS:}
		\begin{itemize}
			\item Lawsuits and liability
			\item Fines and penalties
			\item Reputational damage
		\end{itemize}
		
		\vspace{0.1cm}
		\textbf{HARMS TO SOCIETY:}
		\begin{itemize}
			\item National security risks
			\item Democratic threats
			\item Blackmail of officials
		\end{itemize}
		
		\textcolor{myred}{\textbf{EXAM TRAP:}} Privacy violations cause \textcolor{myblue}{wrongs and harms} to individuals, institutions, and society!
	\end{frame}
	
	\begin{frame}{T54: Personal Data as Toxic Asset}
		\fontsize{6.5}{7.5}\selectfont
		\textbf{TOPIC: Personal Data - A Toxic Asset}
		
		\vspace{0.1cm}
		\textbf{WHY "TOXIC"?}
		\begin{enumerate}
			\item \textcolor{myblue}{Cheap to Mine:}
			\begin{itemize}
				\item Easy to collect
				\item Low cost
				\item Widespread availability
			\end{itemize}
			\item \textcolor{myblue}{Valuable:}
			\begin{itemize}
				\item Profitable for companies
				\item Marketing and targeting
				\item Business insights
			\end{itemize}
			\item \textcolor{myblue}{Sensitive:}
			\begin{itemize}
				\item Privacy concerns
				\item Vulnerable to abuse
				\item High risk
			\end{itemize}
			\item \textcolor{myblue}{Hard to Keep Safe:}
			\begin{itemize}
				\item Security challenges
				\item Attacker advantage
				\item Constant threats
			\end{itemize}
		\end{enumerate}
		
		\vspace{0.1cm}
		\textbf{KEY INSIGHT:}
		\begin{itemize}
			\item Defenders are at disadvantage
			\item Attacker chooses time and method
			\item With enough resources, breach is inevitable
		\end{itemize}
		
		\textcolor{myred}{\textbf{EXAM TRAP:}} Personal data is a \textcolor{myblue}{toxic asset} - valuable but risky!
	\end{frame}
	
	\begin{frame}{T55: Data Minimization}
		\fontsize{6.5}{7.5}\selectfont
		\textbf{TOPIC: The Principle of Data Minimization}
		
		\vspace{0.1cm}
		\textbf{CORE IDEA:}
		\begin{itemize}
			\item Collect only necessary data
			\item Fulfill specific purpose
			\item Most effective privacy protection
		\end{itemize}
		
		\vspace{0.1cm}
		\textbf{WHY IT MATTERS:}
		\begin{enumerate}
			\item \textcolor{myblue}{Risk Reduction:}
			\begin{itemize}
				\item Less data = less risk
				\item Fewer targets
				\item Reduced exposure
			\end{itemize}
			\item \textcolor{myblue}{Respect for Persons:}
			\begin{itemize}
				\item Privacy rights
				\item Avoid unnecessary intrusion
				\item Proportionality
			\end{itemize}
			\item \textcolor{myblue}{Regulatory Compliance:}
			\begin{itemize}
				\item GDPR requirement
				\item Privacy laws
				\item Best practice
			\end{itemize}
		\end{enumerate}
		
		\vspace{0.1cm}
		\textbf{KEY POINT:}
		\begin{itemize}
			\item Only collect data when benefit outweighs disadvantages
			\item Data subjects deserve protection
			\item Minimal data collection is respectful
		\end{itemize}
		
		\textcolor{myred}{\textbf{EXAM TRAP:}} Data minimization = \textcolor{myblue}{collect only necessary data}!
	\end{frame}
	
	\begin{frame}{T56: Right to be Forgotten}
		\fontsize{6.5}{7.5}\selectfont
		\textbf{TOPIC: The Right to be Forgotten}
		
		\vspace{0.1cm}
		\textbf{DEFINITION:}
		\begin{itemize}
			\item Person's right to have private information removed
			\item From Internet searches and other directories
			\item When data is inadequate, irrelevant, or excessive
		\end{itemize}
		
		\vspace{0.1cm}
		\textbf{WHEN IT APPLIES:}
		\begin{enumerate}
			\item \textcolor{myblue}{Incorrect Information:}
			\begin{itemize}
				\item Data that is wrong
				\item Inaccurate records
				\item False information
			\end{itemize}
			\item \textcolor{myblue}{Irrelevant Information:}
			\begin{itemize}
				\item No longer relevant
				\item Outdated
				\item No legitimate purpose
			\end{itemize}
			\item \textcolor{myblue}{Excessive Information:}
			\begin{itemize}
				\item More than needed
				\item Proportionality
				\item Over-collection
			\end{itemize}
		\end{enumerate}
		
		\textcolor{myred}{\textbf{EXAM TRAP:}} Right to be forgotten = \textcolor{myblue}{removal of private information when inadequate or irrelevant}!
	\end{frame}
	
	\begin{frame}{T57: Control Over Data}
		\fontsize{6.5}{7.5}\selectfont
		\textbf{TOPIC: Giving People Control Over Their Data}
		
		\vspace{0.1cm}
		\textbf{KEY RIGHTS:}
		\begin{enumerate}
			\item \textcolor{myblue}{Consent Management:}
			\begin{itemize}
				\item Easy to deny consent
				\item Easy to withdraw consent
				\item Clear options
			\end{itemize}
			\item \textcolor{myblue}{Access:}
			\begin{itemize}
				\item See personal data
				\item See inferences made
				\item Understand data use
			\end{itemize}
			\item \textcolor{myblue}{Transfer:}
			\begin{itemize}
				\item Data portability
				\item Move data to other services
				\item Freedom to choose
			\end{itemize}
			\item \textcolor{myblue}{Deletion:}
			\begin{itemize}
				\item Right to delete data
				\item Remove personal information
				\item Erasure
			\end{itemize}
		\end{enumerate}
		
		\textcolor{myred}{\textbf{EXAM TRAP:}} Control over data = \textcolor{myblue}{consent, access, transfer, and deletion rights}!
	\end{frame}
	
	\begin{frame}{T58: Contextual Integrity}
		\fontsize{6.5}{7.5}\selectfont
		\textbf{TOPIC: Contextual Integrity in Privacy}
		
		\vspace{0.1cm}
		\textbf{CORE IDEA:}
		\begin{itemize}
			\item Data use should adhere to contextual norms
			\item When data is transferred to different context, norms are violated
			\item Context matters
		\end{itemize}
		
		\vspace{0.1cm}
		\textbf{EXAMPLES:}
		\begin{enumerate}
			\item \textcolor{myblue}{Banking Context:}
			\begin{itemize}
				\item Data given for transaction
				\item Should not be sold to marketing company
				\item Violates expectations
			\end{itemize}
			\item \textcolor{myblue}{Healthcare Context:}
			\begin{itemize}
				\item Patient data given to doctor
				\item Should not end up with data brokers
				\item Violates trust
			\end{itemize}
			\item \textcolor{myblue}{Social Media Context:}
			\begin{itemize}
				\item Expectations about data use
				\item Context-specific norms
				\item Violations cause backlash
			\end{itemize}
		\end{enumerate}
		
		\textcolor{myred}{\textbf{EXAM TRAP:}} Contextual integrity = \textcolor{myblue}{data use should respect context-specific norms}!
	\end{frame}
	
	\begin{frame}{T59: Data Deletion Practices}
		\fontsize{6.5}{7.5}\selectfont
		\textbf{TOPIC: Best Practices for Data Deletion}
		
		\vspace{0.1cm}
		\textbf{WHY DELETE DATA:}
		\begin{enumerate}
			\item \textcolor{myblue}{Protect Individuals:}
			\begin{itemize}
				\item Reduce exposure
				\item Privacy protection
				\item Security
			\end{itemize}
			\item \textcolor{myblue}{Maintain Accuracy:}
			\begin{itemize}
				\item Personal data changes quickly
				\item People change jobs, houses, preferences
				\item Outdated data is inaccurate
			\end{itemize}
			\item \textcolor{myblue}{Reduce Liability:}
			\begin{itemize}
				\item Less data = less risk
				\item Fewer targets
				\item Smaller breach impact
			\end{itemize}
		\end{enumerate}
		
		\vspace{0.1cm}
		\textbf{BEST PRACTICE:}
		\begin{itemize}
			\item Regular deletion schedules
			\item Expiry dates for data
			\item Not collected with intention of keeping forever
			\item Routine data deletion
		\end{itemize}
		
		\textcolor{myred}{\textbf{EXAM TRAP:}} Delete data \textcolor{myblue}{as soon as no longer necessary} - routine practice!
	\end{frame}
	
	\begin{frame}{T60: Privacy Principles Summary}
		\fontsize{6.5}{7.5}\selectfont
		\textbf{TOPIC: Summary of Privacy Principles}
		
		\vspace{0.1cm}
		\textbf{KEY PRINCIPLES:}
		\begin{enumerate}
			\item \textcolor{myblue}{Right to Privacy:}
			\begin{itemize}
				\item Moral and legal right
				\item Universal Declaration of Human Rights
				\item Constitutional protections
			\end{itemize}
			\item \textcolor{myblue}{Data Minimization:}
			\begin{itemize}
				\item Collect only necessary data
				\item Most effective protection
			\end{itemize}
			\item \textcolor{myblue}{Right to be Forgotten:}
			\begin{itemize}
				\item Remove inadequate data
				\item Irrelevant or excessive
			\end{itemize}
			\item \textcolor{myblue}{Control Over Data:}
			\begin{itemize}
				\item Consent, access, transfer, deletion
			\end{itemize}
			\item \textcolor{myblue}{Contextual Integrity:}
			\begin{itemize}
				\item Respect context-specific norms
			\end{itemize}
			\item \textcolor{myblue}{Data Deletion:}
			\begin{itemize}
				\item Delete when no longer necessary
			\end{itemize}
			\item \textcolor{myblue}{Data Security:}
			\begin{itemize}
				\item Encryption, anonymization, differential privacy
			\end{itemize}
		\end{enumerate}
		
		\textcolor{myred}{\textbf{EXAM TRAP:}} Multiple principles guide \textcolor{myblue}{ethical privacy practices}!
	\end{frame}
	
	% ============================================================
	% SECTION 9: GOVERNANCE CHALLENGES
	% ============================================================
	
	\begin{frame}{T61: AI Governance Challenges}
		\fontsize{6.5}{7.5}\selectfont
		\textbf{TOPIC: Why AI is Difficult to Govern}
		
		\vspace{0.1cm}
		\textbf{KEY CHALLENGES:}
		\begin{enumerate}
			\item \textcolor{myblue}{Power Asymmetries:}
			\begin{itemize}
				\item Tech companies more powerful than governments
				\item Wealth, influence, expertise
				\item Regulation difficulties
			\end{itemize}
			\item \textcolor{myblue}{Lack of Transparency:}
			\begin{itemize}
				\item Private company development
				\item Limited public information
				\item Safety testing unknown
			\end{itemize}
			\item \textcolor{myblue}{Lack of Interpretability:}
			\begin{itemize}
				\item Neural networks are "black boxes"
				\item Proxy variables obscure
				\item Hard to audit
			\end{itemize}
			\item \textcolor{myblue}{Lack of Ethics Structures:}
			\begin{itemize}
				\item No required ethics education
				\item No licensing
				\item No oversight committees
			\end{itemize}
		\end{enumerate}
		
		\textcolor{myred}{\textbf{EXAM TRAP:}} AI governance faces \textcolor{myblue}{multiple interconnected challenges}!
	\end{frame}
	
	\begin{frame}{T62: Power Asymmetries}
		\fontsize{6.5}{7.5}\selectfont
		\textbf{TOPIC: Power Asymmetries in AI Governance}
		
		\vspace{0.1cm}
		\textbf{THE PROBLEM:}
		\begin{itemize}
			\item Tech companies often more powerful than governments
			\item Wealth: Massive resources
			\item Influence: Lobbying and public relations
			\item Expertise: Advanced technical knowledge
		\end{itemize}
		
		\vspace{0.1cm}
		\textbf{CONSEQUENCES:}
		\begin{enumerate}
			\item \textcolor{myblue}{Regulation Challenges:}
			\begin{itemize}
				\item Governments struggle to keep up
				\item Companies resist regulation
				\item Asymmetric negotiation
			\end{itemize}
			\item \textcolor{myblue}{Information Asymmetry:}
			\begin{itemize}
				\item Companies know more
				\item Governments lack technical expertise
				\item Hard to regulate effectively
			\end{itemize}
			\item \textcolor{myblue}{Accountability Gaps:}
			\begin{itemize}
				\item Who holds companies accountable?
				\item Regulatory capture risk
				\item Limited enforcement
			\end{itemize}
		\end{enumerate}
		
		\textcolor{myred}{\textbf{EXAM TRAP:}} Tech companies may be \textcolor{myblue}{more powerful than some governments}!
	\end{frame}
	
	\begin{frame}{T63: Lack of Transparency}
		\fontsize{6.5}{7.5}\selectfont
		\textbf{TOPIC: Transparency Challenges in AI}
		
		\vspace{0.1cm}
		\textbf{WHY IT'S A PROBLEM:}
		\begin{itemize}
			\item Most AI systems developed by private companies
			\item Not subject to transparency requirements
			\item Limited information about:
			\begin{itemize}
				\item Training datasets
				\item Testing and tuning
				\item Safety measures
			\end{itemize}
		\end{itemize}
		
		\vspace{0.1cm}
		\textbf{CONSEQUENCES:}
		\begin{enumerate}
			\item \textcolor{myblue}{No Public Oversight:}
			\begin{itemize}
				\item Academics can't study systems
				\item Journalists can't investigate
				\item Public is in the dark
			\end{itemize}
			\item \textcolor{myblue}{Regulatory Challenges:}
			\begin{itemize}
				\item Regulators lack information
				\item Can't assess risks
				\item Enforcement difficulties
			\end{itemize}
			\item \textcolor{myblue}{Trust Issues:}
			\begin{itemize}
				\item Lack of transparency erodes trust
				\item Suspicion about practices
				\item Public concern
			\end{itemize}
		\end{enumerate}
		
		\textcolor{myred}{\textbf{EXAM TRAP:}} Private companies are \textcolor{myblue}{not subject to same transparency requirements}!
	\end{frame}
	
	\begin{frame}{T64: Lack of Interpretability}
		\fontsize{6.5}{7.5}\selectfont
		\textbf{TOPIC: Interpretability Challenges}
		
		\vspace{0.1cm}
		\textbf{THE BLACK-BOX PROBLEM:}
		\begin{itemize}
			\item Neural networks are "black boxes"
			\item Even experts can't fully understand
			\item Backpropagation doesn't provide insight
			\item Hard to trace decisions
		\end{itemize}
		
		\vspace{0.1cm}
		\textbf{WHY IT MATTERS:}
		\begin{enumerate}
			\item \textcolor{myblue}{Accountability:}
			\begin{itemize}
				\item Can't hold accountable what you can't understand
				\item Responsibility gaps
				\item Who is responsible?
			\end{itemize}
			\item \textcolor{myblue}{Auditing:}
			\begin{itemize}
				\item Hard to audit black boxes
				\item New methods needed
				\item Limited effectiveness
			\end{itemize}
			\item \textcolor{myblue}{Governance:}
			\begin{itemize}
				\item Can't govern what you can't understand
				\item Regulatory challenges
				\item Oversight difficulties
			\end{itemize}
		\end{enumerate}
		
		\textcolor{myred}{\textbf{EXAM TRAP:}} Neural networks are \textcolor{myblue}{black boxes} - even experts can't fully understand them!
	\end{frame}
	
	\begin{frame}{T65: Proxy Variables Challenge}
		\fontsize{6.5}{7.5}\selectfont
		\textbf{TOPIC: Proxy Variables and Governance}
		
		\vspace{0.1cm}
		\textbf{WHAT ARE PROXIES?}
		\begin{itemize}
			\item Techniques to simplify training
			\item Represent complex objectives with simpler metrics
			\item Example: Word count as proxy for article quality
		\end{itemize}
		
		\vspace{0.1cm}
		\textbf{THE CHALLENGE:}
		\begin{enumerate}
			\item \textcolor{myblue}{Alignment Issues:}
			\begin{itemize}
				\item Proxy may not align with true objective
				\item Optimization for wrong goal
				\item Unintended consequences
			\end{itemize}
			\item \textcolor{myblue}{Opacity:}
			\begin{itemize}
				\item If proxies are unknown
				\item Hard to understand model behavior
				\item Governance difficulties
			\end{itemize}
			\item \textcolor{myblue}{Scale and Complexity:}
			\begin{itemize}
				\item Billions of parameters
				\item Hard to trace processes
				\item Audit challenges
			\end{itemize}
		\end{enumerate}
		
		\textcolor{myred}{\textbf{EXAM TRAP:}} Unknown proxies \textcolor{myblue}{make governance difficult}!
	\end{frame}
	
	\begin{frame}{T66: Lack of Ethics Structures}
		\fontsize{6.5}{7.5}\selectfont
		\textbf{TOPIC: Lack of AI Ethics Structures}
		
		\vspace{0.1cm}
		\textbf{COMPARISON WITH MEDICAL ETHICS:}
		\begin{tabular}{|p{4cm}|p{4cm}|}
			\hline
			\textbf{Medical Ethics} & \textbf{AI Ethics} \\
			\hline
			Bioethicists & Few AI ethicists \\
			Ethical codes & Developing codes \\
			Ethics committees & Rare committees \\
			Required coursework & Optional courses \\
			Licensed professionals & No licensing \\
			International codes & Fragmented guidance \\
			\hline
		\end{tabular}
		
		\vspace{0.1cm}
		\textbf{KEY GAP:}
		\begin{itemize}
			\item Computer scientists can graduate without ethics training
			\item No certification or licensing
			\item Boards rarely include AI ethicists
			\item No oversight committees
		\end{itemize}
		
		\textcolor{myred}{\textbf{EXAM TRAP:}} AI ethics lacks the \textcolor{myblue}{established structures} of medical ethics!
	\end{frame}
	
	\begin{frame}{T67: Lack of Regulation}
		\fontsize{6.5}{7.5}\selectfont
		\textbf{TOPIC: The Regulatory Gap}
		
		\vspace{0.1cm}
		\textbf{CURRENT STATE:}
		\begin{itemize}
			\item Few national laws (China exception)
			\item No specialized agencies
			\item No government-mandated auditing
			\item Voluntary frameworks only
		\end{itemize}
		
		\vspace{0.1cm}
		\textbf{HISTORICAL PARALLEL:}
		\begin{itemize}
			\item FDA established for food and drugs
			\item Other regulatory agencies developed
			\item AI may follow similar path
			\item EU leading with AI Office
		\end{itemize}
		
		\vspace{0.1cm}
		\textbf{WHY REGULATION MATTERS:}
		\begin{itemize}
			\item AI is international technology
			\item Cross-border data flows
			\item Need for cooperation
			\item International standards needed
		\end{itemize}
		
		\textcolor{myred}{\textbf{EXAM TRAP:}} Few national AI laws exist \textcolor{myblue}{(except China)} - no specialized agencies!
	\end{frame}
	
	\begin{frame}{T68: Unpredictability}
		\fontsize{6.5}{7.5}\selectfont
		\textbf{TOPIC: Unpredictability in AI Systems}
		
		\vspace{0.1cm}
		\textbf{EMERGENT BEHAVIOR:}
		\begin{itemize}
			\item Properties not explicitly coded
			\item Capabilities that surprise
			\item Mistakes that are unexpected
			\item Interactions that are unforeseen
		\end{itemize}
		
		\vspace{0.1cm}
		\textbf{CONSEQUENCES:}
		\begin{enumerate}
			\item \textcolor{myblue}{Safety Risks:}
			\begin{itemize}
				\item Unintended outcomes
				\item Harmful behavior
				\item System failures
			\end{itemize}
			\item \textcolor{myblue}{Governance Challenges:}
			\begin{itemize}
				\item Hard to regulate unpredictable systems
				\item Misuse difficult to prevent
				\item Unknown risks
			\end{itemize}
			\item \textcolor{myblue}{Trust Issues:}
			\begin{itemize}
				\item Surprises undermine trust
				\item Unpredictable = unreliable
				\item Adoption barriers
			\end{itemize}
		\end{enumerate}
		
		\textcolor{myred}{\textbf{EXAM TRAP:}} Emergent behavior = \textcolor{myblue}{unexpected capabilities not explicitly coded}!
	\end{frame}
	
	\begin{frame}{T69: Truth-Tracking Limitations}
		\fontsize{6.5}{7.5}\selectfont
		\textbf{TOPIC: LLMs and Truth-Tracking}
		
		\vspace{0.1cm}
		\textbf{HOW LLMS WORK:}
		\begin{itemize}
			\item Statistical inferences and probabilistic guesses
			\item Designed for plausible responses
			\item Not based on truth or knowledge
		\end{itemize}
		
		\vspace{0.1cm}
		\textbf{RISKS:}
		\begin{enumerate}
			\item \textcolor{myblue}{Misinformation:}
			\begin{itemize}
				\item Plausible but false information
				\item Hard to detect
				\item Spread at scale
			\end{itemize}
			\item \textcolor{myblue}{Physical Safety:}
			\begin{itemize}
				\item Incorrect medical advice
				\item Dangerous recommendations
				\item Real-world harm
			\end{itemize}
			\item \textcolor{myblue}{Legal Risks:}
			\begin{itemize}
				\item Libel and defamation
				\item Copyright issues
				\item Speech-related risks
			\end{itemize}
		\end{enumerate}
		
		\textcolor{myred}{\textbf{EXAM TRAP:}} LLMs give \textcolor{myblue}{plausible, not necessarily truthful} responses!
	\end{frame}
	
	\begin{frame}{T70: Privacy and Copyright Governance}
		\fontsize{6.5}{7.5}\selectfont
		\textbf{TOPIC: Privacy and Copyright Challenges}
		
		\vspace{0.1cm}
		\textbf{PRIVACY CONCERNS:}
		\begin{itemize}
			\item Data scraped from internet
			\item Personal data of unsuspecting users
			\item GDPR compliance uncertain
			\item Right to access, modify, delete unclear
		\end{itemize}
		
		\vspace{0.1cm}
		\textbf{COPYRIGHT CONCERNS:}
		\begin{itemize}
			\item LLMs ingest copyrighted material
			\item Books, articles, images
			\item Lawsuits in process
			\item Legal uncertainty
		\end{itemize}
		
		\vspace{0.1cm}
		\textbf{GOVERNANCE NIGHTMARE:}
		\begin{itemize}
			\item Machine-created speech
			\item Company ownership
			\item Data acquisition unclear
			\item Regulatory challenges
		\end{itemize}
		
		\textcolor{myred}{\textbf{EXAM TRAP:}} Scraping data raises \textcolor{myblue}{privacy and copyright concerns}!
	\end{frame}
	
	% ============================================================
	% SECTION 10: REGULATORY LANDSCAPE
	% ============================================================
	
	\begin{frame}{T71: Regulatory Overview}
		\fontsize{6.5}{7.5}\selectfont
		\textbf{TOPIC: Global AI Regulatory Landscape}
		
		\vspace{0.1cm}
		\textbf{KEY PLAYERS:}
		\begin{enumerate}
			\item \textcolor{myblue}{European Union:}
			\begin{itemize}
				\item GDPR (Privacy)
				\item DMA (Digital Markets)
				\item DSA (Digital Services)
				\item AI Act (Comprehensive AI Law)
			\end{itemize}
			\item \textcolor{myblue}{United States:}
			\begin{itemize}
				\item No federal privacy law
				\item State-level privacy laws
				\item AI Executive Order
				\item NIST AI RMF
			\end{itemize}
			\item \textcolor{myblue}{China:}
			\begin{itemize}
				\item Generative AI Measures
				\item PIPL (Privacy)
				\item Algorithmic Recommendation Rules
			\end{itemize}
		\end{enumerate}
		
		\textcolor{myred}{\textbf{EXAM TRAP:}} Different regions have \textcolor{myblue}{different regulatory approaches}!
	\end{frame}
	
	\begin{frame}{T72: GDPR Overview}
		\fontsize{6.5}{7.5}\selectfont
		\textbf{TOPIC: European GDPR}
		
		\vspace{0.1cm}
		\textbf{KEY FACTS:}
		\begin{itemize}
			\item Implemented May 2018
			\item Regulates personal data
			\item Applies to data controllers and processors
			\item Extraterritorial jurisdiction
		\end{itemize}
		
		\vspace{0.1cm}
		\textbf{DATA SUBJECT RIGHTS:}
		\begin{enumerate}
			\item Receive concise information
			\item Access personal data
			\item Request erasure
			\item Object to processing
		\end{enumerate}
		
		\vspace{0.1cm}
		\textbf{OBLIGATIONS:}
		\begin{itemize}
			\item Notify breaches within 72 hours
			\item "Legitimate interest" required
			\item Significant fines for non-compliance
		\end{itemize}
		
		\textcolor{myred}{\textbf{EXAM TRAP:}} GDPR = \textcolor{myblue}{European data protection regulation} with global reach!
	\end{frame}
	
	\begin{frame}{T73: DMA Overview}
		\fontsize{6.5}{7.5}\selectfont
		\textbf{TOPIC: Digital Markets Act (DMA)}
		
		\vspace{0.1cm}
		\textbf{KEY FACTS:}
		\begin{itemize}
			\item Passed by European Union in 2022
			\item Aims for fair and competitive digital marketspace
			\item Targets "gatekeeper" platforms
		\end{itemize}
		
		\vspace{0.1cm}
		\textbf{KEY PROVISIONS:}
		\begin{enumerate}
			\item \textcolor{myblue}{Ban on Self-Preferencing:}
			\begin{itemize}
				\item Can't favor own services
				\item No ranking own products higher
				\item No making switching difficult
			\end{itemize}
			\item \textcolor{myblue}{Open Access to Data:}
			\begin{itemize}
				\item Provide data to third parties
				\item Enable innovative services
				\item Foster competition
			\end{itemize}
			\item \textcolor{myblue}{Transparency Obligations:}
			\begin{itemize}
				\item Transparent about algorithms
				\item Allow understanding of operations
			\end{itemize}
			\item \textcolor{myblue}{No Tying and Bundling:}
			\begin{itemize}
				\item Can't require unwanted services
				\item Core services accessible separately
			\end{itemize}
		\end{enumerate}
		
		\textcolor{myred}{\textbf{EXAM TRAP:}} DMA targets \textcolor{myblue}{gatekeeper platforms} for fair competition!
	\end{frame}
	
	\begin{frame}{T74: DSA Overview}
		\fontsize{6.5}{7.5}\selectfont
		\textbf{TOPIC: Digital Services Act (DSA)}
		
		\vspace{0.1cm}
		\textbf{KEY FACTS:}
		\begin{itemize}
			\item Adopted October 2022
			\item Came into force 2023
			\item Applies to online intermediaries and platforms
			\item 19 companies currently beholden
		\end{itemize}
		
		\vspace{0.1cm}
		\textbf{KEY OBLIGATIONS:}
		\begin{enumerate}
			\item \textcolor{myblue}{Prevent Illegal Content:}
			\begin{itemize}
				\item Proactively identify and remove
				\item Hate speech, CSAM, counterfeit goods
				\item Clear reporting mechanisms
			\end{itemize}
			\item \textcolor{myblue}{Transparency:}
			\begin{itemize}
				\item Publish content-moderation policies
				\item Provide access to moderation data
			\end{itemize}
			\item \textcolor{myblue}{Address Disinformation:}
			\begin{itemize}
				\item Labels on political advertising
				\item Promote fact-checking
			\end{itemize}
			\item \textcolor{myblue}{Protect from Algorithmic Bias:}
			\begin{itemize}
				\item Ensure algorithms not biased
				\item Transparent about algorithms
			\end{itemize}
			\item \textcolor{myblue}{Opt-Out of Personalized Content:}
			\begin{itemize}
				\item Very large platforms must allow
				\item No tracking for recommendations
			\end{itemize}
		\end{enumerate}
		
		\textcolor{myred}{\textbf{EXAM TRAP:}} DSA addresses \textcolor{myblue}{illegal content, transparency, and disinformation}!
	\end{frame}
	
	\begin{frame}{T75: EU AI Act Overview}
		\fontsize{6.5}{7.5}\selectfont
		\textbf{TOPIC: The EU AI Act}
		
		\vspace{0.1cm}
		\textbf{KEY FACTS:}
		\begin{itemize}
			\item Approved March 2024
			\item Council approved May 2024
			\item Establishes common regulatory framework
			\item Risk-based approach
		\end{itemize}
		
		\vspace{0.1cm}
		\textbf{RISK CATEGORIES:}
		\begin{enumerate}
			\item \textcolor{myblue}{Unacceptable Risk (Banned):}
			\begin{itemize}
				\item Biometric categorization using sensitive characteristics
				\item Untargeted facial image scraping
				\item Emotion recognition in workplace
				\item Social scoring
				\item Subliminal manipulation
			\end{itemize}
			\item \textcolor{myblue}{High-Risk:}
			\begin{itemize}
				\item Impact assessments required
				\item Thorough records
				\item Oversight measures
			\end{itemize}
			\item \textcolor{myblue}{Limited/Minimal Risk:}
			\begin{itemize}
				\item Transparency obligations
				\item Labeling requirements
			\end{itemize}
		\end{enumerate}
		
		\textcolor{myred}{\textbf{EXAM TRAP:}} EU AI Act = \textcolor{myblue}{first comprehensive AI regulation} - risk-based approach!
	\end{frame}
	
	\begin{frame}{T76: EU AI Act Requirements}
		\fontsize{6.5}{7.5}\selectfont
		\textbf{TOPIC: EU AI Act - Specific Requirements}
		
		\vspace{0.1cm}
		\textbf{TRANSPARENCY REQUIREMENTS:}
		\begin{itemize}
			\item Notify people when interacting with chatbot
			\item Notify when using biometric systems
			\item Notify when using emotion recognition
			\item Label deepfakes and AI-generated content
			\item Design systems to make AI-generated media detectable
		\end{itemize}
		
		\vspace{0.1cm}
		\textbf{HIGH-RISK REQUIREMENTS:}
		\begin{itemize}
			\item Impact assessment on rights
			\item Keep thorough records of:
			\begin{itemize}
				\item Datasets used
				\item Programming and training methodologies
				\item Oversight measures
			\end{itemize}
		\end{itemize}
		
		\vspace{0.1cm}
		\textbf{PENALTIES:}
		\begin{itemize}
			\item €35 million or 7\% of global turnover
			\item €7.5 million or 1.5\% of turnover
			\item Depends on offense and company size
		\end{itemize}
		
		\textcolor{myred}{\textbf{EXAM TRAP:}} EU AI Act has \textcolor{myblue}{significant penalties} for non-compliance!
	\end{frame}
	
	\begin{frame}{T77: US Privacy Regulations}
		\fontsize{6.5}{7.5}\selectfont
		\textbf{TOPIC: US Privacy Regulations}
		
		\vspace{0.1cm}
		\textbf{NO FEDERAL PRIVACY LAW:}
		\begin{itemize}
			\item United States lacks comprehensive federal privacy law
			\item Patchwork of state-level regulations
			\item HIPAA for health information
		\end{itemize}
		
		\vspace{0.1cm}
		\textbf{KEY STATE LAWS:}
		\begin{enumerate}
			\item \textcolor{myblue}{California Consumer Privacy Act (CCPA):}
			\begin{itemize}
				\item Most comprehensive state privacy law
				\item Right to know, access, delete
				\item Right to opt out of sale
				\item Right to non-discrimination
			\end{itemize}
			\item \textcolor{myblue}{Virginia Consumer Data Protection Act (CDPA):}
			\begin{itemize}
				\item Modeled after CCPA
				\item Additional protections for sensitive data
				\item Genetic and biometric data
			\end{itemize}
			\item \textcolor{myblue}{Colorado Privacy Act (CPA):}
			\begin{itemize}
				\item Right to opt out of targeted advertising
				\item Data minimization provisions
			\end{itemize}
			\item \textcolor{myblue}{Illinois BIPA:}
			\begin{itemize}
				\item Biometric Information Privacy Act
				\item Notification, access, deletion rights
				\item Civil remedies for violations
			\end{itemize}
		\end{enumerate}
		
		\textcolor{myred}{\textbf{EXAM TRAP:}} US has \textcolor{myblue}{patchwork of state privacy laws} - no federal law!
	\end{frame}
	
	\begin{frame}{T78: US AI Regulation}
		\fontsize{6.5}{7.5}\selectfont
		\textbf{TOPIC: US AI Regulatory Framework}
		
		\vspace{0.1cm}
		\textbf{KEY INITIATIVES:}
		\begin{enumerate}
			\item \textcolor{myblue}{Executive Order on AI (2023):}
			\begin{itemize}
				\item Safe, secure, trustworthy development
				\item Address AI-enabled harms
				\item International cooperation
				\item Safety testing requirements
			\end{itemize}
			\item \textcolor{myblue}{NIST AI RMF:}
			\begin{itemize}
				\item Voluntary framework
				\item Govern, map, measure, manage
				\item Flexible and adaptable
				\item Non-prescriptive
			\end{itemize}
			\item \textcolor{myblue}{AI Safety and Security Board:}
			\begin{itemize}
				\item Department of Homeland Security
				\item Critical infrastructure focus
				\item Recommendations and guidance
			\end{itemize}
		\end{enumerate}
		
		\textcolor{myred}{\textbf{EXAM TRAP:}} US has \textcolor{myblue}{executive orders and voluntary frameworks} - no federal AI law yet!
	\end{frame}
	
	\begin{frame}{T79: China AI Regulation}
		\fontsize{6.5}{7.5}\selectfont
		\textbf{TOPIC: China AI Regulatory Framework}
		
		\vspace{0.1cm}
		\textbf{KEY REGULATIONS:}
		\begin{enumerate}
			\item \textcolor{myblue}{Generative AI Measures (2023):}
			\begin{itemize}
				\item Applies to generative AI services
				\item Prohibits illegal content
				\item Prevents discriminatory content
				\item Protects privacy and personal information
			\end{itemize}
			\item \textcolor{myblue}{Personal Information Protection Law (PIPL):}
			\begin{itemize}
				\item Comprehensive privacy law
				\item Data minimization requirements
				\item User consent requirements
			\end{itemize}
			\item \textcolor{myblue}{Algorithmic Recommendation Rules:}
			\begin{itemize}
				\item Transparency in algorithm decision-making
				\item User consent requirements
				\item Data minimization
			\end{itemize}
		\end{enumerate}
		
		\textcolor{myred}{\textbf{EXAM TRAP:}} China has \textcolor{myblue}{comprehensive AI regulations} - among first globally!
	\end{frame}
	
	\begin{frame}{T80: Regulatory Comparison}
		\fontsize{6.5}{7.5}\selectfont
		\textbf{TOPIC: Comparing Global AI Regulation}
		
		\vspace{0.1cm}
		\textbf{EUROPEAN UNION:}
		\begin{itemize}
			\item Comprehensive, risk-based approach
			\item GDPR for privacy
			\item AI Act for AI specifically
			\item Strong enforcement
			\item Extraterritorial reach
		\end{itemize}
		
		\vspace{0.1cm}
		\textbf{UNITED STATES:}
		\begin{itemize}
			\item No federal AI law
			\item State-level privacy laws
			\item Executive orders and frameworks
			\item Voluntary compliance
			\item Sector-specific approach
		\end{itemize}
		
		\vspace{0.1cm}
		\textbf{CHINA:}
		\begin{itemize}
			\item Comprehensive AI regulations
			\item Generative AI measures
			\item Strong state control
			\item Rapid development
		\end{itemize}
		
		\textcolor{myred}{\textbf{EXAM TRAP:}} EU leads with \textcolor{myblue}{comprehensive regulation}; US relies on \textcolor{myblue}{voluntary frameworks}!
	\end{frame}
	
	% ============================================================
	% SECTION 11: ETHICS AND LAW RELATIONSHIP
	% ============================================================
	
	\begin{frame}{T81: Ethics as Complement to Law}
		\fontsize{6.5}{7.5}\selectfont
		\textbf{TOPIC: How Ethics Complements Law}
		
		\vspace{0.1cm}
		\textbf{ETHICS GROUNDS LAW:}
		\begin{itemize}
			\item Ethics helps ground, inform, and shape law
			\item Societies regulate behavior based on moral acceptability
			\item Ethics helps distinguish between just and unjust laws
		\end{itemize}
		
		\vspace{0.1cm}
		\textbf{LAWS ARE MINIMAL:}
		\begin{itemize}
			\item Establish minimal behavioral requirements
			\item Enable orderly interactions
			\item Provide framework for society
		\end{itemize}
		
		\vspace{0.1cm}
		\textbf{ETHICS GOES BEYOND:}
		\begin{itemize}
			\item Identifies moral issues
			\item Reflects on kind of society we want
			\item Makes recommendations for a good life
			\item More ambitious than law
		\end{itemize}
		
		\vspace{0.1cm}
		\textbf{EXAMPLE:}
		\begin{itemize}
			\item Law: Not illegal to abandon friends
			\item Ethics: Abandoning friends is immoral
			\item Laws provide minimums; ethics provides ideals
		\end{itemize}
		
		\textcolor{myred}{\textbf{EXAM TRAP:}} Ethics \textcolor{myblue}{complements and informs} law - goes beyond minimum requirements!
	\end{frame}
	
	\begin{frame}{T82: Ethics and AI Regulation}
		\fontsize{6.5}{7.5}\selectfont
		\textbf{TOPIC: Ethics in AI Regulation}
		
		\vspace{0.1cm}
		\textbf{RELATIONSHIP:}
		\begin{itemize}
			\item AI ethics gaining importance
			\item Some feel laws also needed
			\item Ethics informs regulatory approaches
		\end{itemize}
		
		\vspace{0.1cm}
		\textbf{EXAMPLES:}
		\begin{itemize}
			\item GDPR: Not designed for AI but relevant
			\item EU AI Act: Based on ethical principles
			\item NIST AI RMF: Incorporates ethical considerations
		\end{itemize}
		
		\vspace{0.1cm}
		\textbf{KEY INSIGHT:}
		\begin{itemize}
			\item Regulations increasingly require ethical considerations
			\item Ethics frameworks underlay responsible AI approaches
			\item Ethical principles guide regulatory development
		\end{itemize}
		
		\textcolor{myred}{\textbf{EXAM TRAP:}} Ethics \textcolor{myblue}{informs and shapes} AI regulation!
	\end{frame}
	
	% ============================================================
	% SECTION 12: PUTTING IT ALL TOGETHER
	% ============================================================
	
	\begin{frame}{T83: Responsible AI Framework}
		\fontsize{6.5}{7.5}\selectfont
		\textbf{TOPIC: Comprehensive Responsible AI Framework}
		
		\vspace{0.1cm}
		\textbf{KEY COMPONENTS:}
		\begin{enumerate}
			\item \textcolor{myblue}{Ethical Foundations:}
			\begin{itemize}
				\item Consequentialism
				\item Deontology
				\item Virtue ethics
			\end{itemize}
			\item \textcolor{myblue}{Core Principles:}
			\begin{itemize}
				\item Nonmaleficence, beneficence, justice, autonomy, explainability
			\end{itemize}
			\item \textcolor{myblue}{Bias Mitigation:}
			\begin{itemize}
				\item Identify sources
				\item Implement strategies
				\item Continuous monitoring
			\end{itemize}
			\item \textcolor{myblue}{Privacy Protection:}
			\begin{itemize}
				\item Data minimization
				\item Security measures
				\item User control
			\end{itemize}
			\item \textcolor{myblue}{Governance:}
			\begin{itemize}
				\item Ethics committees
				\item Auditing
				\item Accountability
			\end{itemize}
		\end{enumerate}
		
		\textcolor{myred}{\textbf{EXAM TRAP:}} Responsible AI requires \textcolor{myblue}{comprehensive approach} across all areas!
	\end{frame}
	
	\begin{frame}{T84: Role of Ethics Committees}
		\fontsize{6.5}{7.5}\selectfont
		\textbf{TOPIC: The Role of Ethics Committees in AI}
		
		\vspace{0.1cm}
		\textbf{PURPOSE:}
		\begin{itemize}
			\item Discuss possible problematic biases
			\item Make decisions about trade-offs
			\item Consider multiple ethical frameworks
			\item Ensure procedural fairness
		\end{itemize}
		
		\vspace{0.1cm}
		\textbf{BENEFITS:}
		\begin{enumerate}
			\item \textcolor{myblue}{Diverse Perspectives:}
			\begin{itemize}
				\item Different viewpoints
				\item Multiple expertise areas
				\item Balanced considerations
			\end{itemize}
			\item \textcolor{myblue}{Structured Decision-Making:}
			\begin{itemize}
				\item Systematic review
				\item Transparent processes
				\item Documented reasoning
			\end{itemize}
			\item \textcolor{myblue}{Accountability:}
			\begin{itemize}
				\item Clear responsibility
				\item Audit trail
				\item Defensible decisions
			\end{itemize}
		\end{enumerate}
		
		\textcolor{myred}{\textbf{EXAM TRAP:}} Ethics committees help \textcolor{myblue}{discuss and address biases} through diverse perspectives!
	\end{frame}
	
	\begin{frame}{T85: Continuous Monitoring}
		\fontsize{6.5}{7.5}\selectfont
		\textbf{TOPIC: Continuous Monitoring and Evaluation}
		
		\vspace{0.1cm}
		\textbf{IMPORTANCE:}
		\begin{itemize}
			\item First strategy to prevent reputational loss
			\item Identify risks proactively
			\item Address them proactively
		\end{itemize}
		
		\vspace{0.1cm}
		\textbf{HOW TO IMPLEMENT:}
		\begin{enumerate}
			\item \textcolor{myblue}{Organizational Governance:}
			\begin{itemize}
				\item Mechanisms for regular monitoring
				\item Clear responsibilities
				\item Reporting structures
			\end{itemize}
			\item \textcolor{myblue}{Risk Assessment:}
			\begin{itemize}
				\item Regular evaluations
				\item Identify new risks
				\item Update mitigation strategies
			\end{itemize}
			\item \textcolor{myblue}{Performance Tracking:}
			\begin{itemize}
				\item Monitor AI system performance
				\item Track fairness metrics
				\item Detect drift or degradation
			\end{itemize}
		\end{enumerate}
		
		\textcolor{myred}{\textbf{EXAM TRAP:}} Continuous monitoring is \textcolor{myblue}{essential for responsible AI}!
	\end{frame}
	
	\begin{frame}{T86: Stakeholder Engagement}
		\fontsize{6.5}{7.5}\selectfont
		\textbf{TOPIC: Stakeholder Engagement and Communication}
		
		\vspace{0.1cm}
		\textbf{DEFINITION:}
		\begin{itemize}
			\item Regularly engage with stakeholders
			\item Customers, employees, regulators
			\item About AI initiatives and their impacts
		\end{itemize}
		
		\vspace{0.1cm}
		\textbf{BENEFITS:}
		\begin{itemize}
			\item Clear communication prevents misunderstandings
			\item Builds trust
			\item Early identification of concerns
			\item Collaborative problem-solving
		\end{itemize}
		
		\vspace{0.1cm}
		\textbf{APPROACH:}
		\begin{itemize}
			\item Regular updates
			\item Transparent about challenges
			\item Responsive to feedback
			\item Inclusive of diverse perspectives
		\end{itemize}
		
		\textcolor{myred}{\textbf{EXAM TRAP:}} Stakeholder engagement \textcolor{myblue}{builds trust and prevents misunderstandings}!
	\end{frame}
	
	\begin{frame}{T87: Crisis Management}
		\fontsize{6.5}{7.5}\selectfont
		\textbf{TOPIC: AI-Related Crisis Management}
		
		\vspace{0.1cm}
		\textbf{DEFINITION:}
		\begin{itemize}
			\item Perform risk analysis
			\item Inform development of crisis-management plan
			\item Tailored to AI-related incidents
		\end{itemize}
		
		\vspace{0.1cm}
		\textbf{COMPONENTS:}
		\begin{itemize}
			\item Communication strategies
			\item Potential steps to rectify issue
			\item Clear lines of responsibility
			\item Escalation procedures
		\end{itemize}
		
		\vspace{0.1cm}
		\textbf{BENEFITS:}
		\begin{itemize}
			\item Faster response to incidents
			\item Minimize reputational damage
			\item Demonstrate competence and responsibility
			\item Rebuild trust more effectively
		\end{itemize}
		
		\textcolor{myred}{\textbf{EXAM TRAP:}} Crisis management plan = \textcolor{myblue}{proactive preparation}!
	\end{frame}
	
	\begin{frame}{T88: Training Leadership}
		\fontsize{6.5}{7.5}\selectfont
		\textbf{TOPIC: Training for Board and C-Suite}
		
		\vspace{0.1cm}
		\textbf{WHY LEADERSHIP EDUCATION MATTERS:}
		\begin{itemize}
			\item Wide range of risks from AI
			\item Financial, legal, reputational consequences
			\item Responsible AI is a strategic imperative
			\item Leadership drives organizational culture
		\end{itemize}
		
		\vspace{0.1cm}
		\textbf{KEY TOPICS FOR TRAINING:}
		\begin{enumerate}
			\item \textcolor{myblue}{AI Risks:}
			\begin{itemize}
				\item Bias and discrimination
				\item Privacy violations
				\item Regulatory non-compliance
				\item Reputational damage
			\end{itemize}
			\item \textcolor{myblue}{Ethical/Responsible AI Practices:}
			\begin{itemize}
				\item Fairness principles
				\item Governance structures
				\item Risk mitigation
			\end{itemize}
			\item \textcolor{myblue}{Strategic Considerations:}
			\begin{itemize}
				\item Competitive advantage
				\item Trust building
				\item Long-term value
			\end{itemize}
		\end{enumerate}
		
		\textcolor{myred}{\textbf{EXAM TRAP:}} Leadership must be \textcolor{myblue}{educated about AI risks and responsible practices}!
	\end{frame}
	
	\begin{frame}{T89: Fairness and Bias Summary}
		\fontsize{6.5}{7.5}\selectfont
		\textbf{TOPIC: Summary - Bias, Discrimination, and Fairness}
		
		\vspace{0.1cm}
		\textbf{KEY CONCEPTS:}
		\begin{enumerate}
			\item \textcolor{myblue}{Bias:}
			\begin{itemize}
				\item Systemic deviation from norm
				\item Can be justifiable or problematic
				\item Sources: Problem, data, modeling, deployment
			\end{itemize}
			\item \textcolor{myblue}{Discrimination:}
			\begin{itemize}
				\item Bias that disfavors protected characteristics
				\item Legally prohibited
				\item Requires proactive measures
			\end{itemize}
			\item \textcolor{myblue}{Fairness:}
			\begin{itemize}
				\item Absence of problematic bias
				\item Group and individual fairness
				\item Cannot be fully automated
			\end{itemize}
		\end{enumerate}
		
		\textcolor{myred}{\textbf{EXAM TRAP:}} Bias management requires \textcolor{myblue}{continuous effort across the AI lifecycle}!
	\end{frame}
	
	\begin{frame}{T90: Bias Mitigation Tools}
		\fontsize{6.5}{7.5}\selectfont
		\textbf{TOPIC: Tools for Bias Mitigation}
		
		\vspace{0.1cm}
		\textbf{AVAILABLE TOOLKITS:}
		\begin{enumerate}
			\item \textcolor{myblue}{Aequitas:}
			\begin{itemize}
				\item Tests models for different bias metrics
				\item Across different groups
				\item Auditing tool
			\end{itemize}
			\item \textcolor{myblue}{AI Fairness 360 (AIF360):}
			\begin{itemize}
				\item IBM's fairness toolkit
				\item Benchmark for testing
				\item Multiple fairness metrics
			\end{itemize}
			\item \textcolor{myblue}{Internal Auditing Systems:}
			\begin{itemize}
				\item Companies can create own systems
				\item Identify potential biases
				\item Continuous monitoring
			\end{itemize}
		\end{enumerate}
		
		\vspace{0.1cm}
		\textbf{EXTERNAL OPTIONS:}
		\begin{itemize}
			\item Third-party auditing services
			\item Independent assessors
			\item Bias reports from providers
			\item Fresh perspective
		\end{itemize}
		
		\textcolor{myred}{\textbf{EXAM TRAP:}} Multiple tools available for \textcolor{myblue}{identifying and mitigating bias}!
	\end{frame}
	
	\begin{frame}{T91: Data Quality Importance}
		\fontsize{6.5}{7.5}\selectfont
		\textbf{TOPIC: Data Quality and Bias Prevention}
		
		\vspace{0.1cm}
		\textbf{KEY DATA QUALITY FACTORS:}
		\begin{enumerate}
			\item \textcolor{myblue}{Diversity:}
			\begin{itemize}
				\item Include all relevant groups
				\item Avoid homogeneity
				\item Represent all populations
			\end{itemize}
			\item \textcolor{myblue}{Accuracy:}
			\begin{itemize}
				\item Correct and reliable data
				\item Avoid errors
				\item Verify sources
			\end{itemize}
			\item \textcolor{myblue}{Representativeness:}
			\begin{itemize}
				\item Reflect target population
				\item Avoid sampling bias
				\item Generalizable results
			\end{itemize}
			\item \textcolor{myblue}{Currency:}
			\begin{itemize}
				\item Up-to-date data
				\item Reflect current reality
				\item Avoid outdated patterns
			\end{itemize}
		\end{enumerate}
		
		\vspace{0.1cm}
		\textbf{SYNTHETIC DATA:}
		\begin{itemize}
			\item Fabricated datasets
			\item No privacy risks
			\item Can be designed to avoid biases
			\item May be less precise
		\end{itemize}
		
		\textcolor{myred}{\textbf{EXAM TRAP:}} Data quality helps \textcolor{myblue}{avoid biases} through diversity and representativeness!
	\end{frame}
	
	\begin{frame}{T92: Algorithmic Auditing}
		\fontsize{6.5}{7.5}\selectfont
		\textbf{TOPIC: The Importance of Algorithmic Auditing}
		
		\vspace{0.1cm}
		\textbf{WHY AUDITING MATTERS:}
		\begin{itemize}
			\item Likely the best way to identify biases
			\item Fresh and diverse perspective
			\item Independent assessment
			\item Continuous improvement
		\end{itemize}
		
		\vspace{0.1cm}
		\textbf{WHAT AUDITING INCLUDES:}
		\begin{enumerate}
			\item \textcolor{myblue}{Technological Tools:}
			\begin{itemize}
				\item Statistical analyses
				\item Fairness metrics
				\item Bias detection
			\end{itemize}
			\item \textcolor{myblue}{Human Review:}
			\begin{itemize}
				\item Diverse perspectives
				\item Contextual understanding
				\item Qualitative assessment
			\end{itemize}
			\item \textcolor{myblue}{Documentation:}
			\begin{itemize}
				\item Findings and recommendations
				\item Remediation plans
				\item Continuous monitoring
			\end{itemize}
		\end{enumerate}
		
		\textcolor{myred}{\textbf{EXAM TRAP:}} Algorithmic auditing is \textcolor{myblue}{best way to identify and correct biases}!
	\end{frame}
	
	\begin{frame}{T93: Privacy Principles Summary}
		\fontsize{6.5}{7.5}\selectfont
		\textbf{TOPIC: Summary of Privacy Principles}
		
		\vspace{0.1cm}
		\textbf{KEY PRINCIPLES:}
		\begin{enumerate}
			\item \textcolor{myblue}{Right to Privacy:}
			\begin{itemize}
				\item Moral and legal right
				\item Universal Declaration of Human Rights
			\end{itemize}
			\item \textcolor{myblue}{Data Minimization:}
			\begin{itemize}
				\item Collect only necessary data
			\end{itemize}
			\item \textcolor{myblue}{Right to be Forgotten:}
			\begin{itemize}
				\item Remove inadequate data
			\end{itemize}
			\item \textcolor{myblue}{Control Over Data:}
			\begin{itemize}
				\item Consent, access, transfer, deletion
			\end{itemize}
			\item \textcolor{myblue}{Contextual Integrity:}
			\begin{itemize}
				\item Respect context-specific norms
			\end{itemize}
			\item \textcolor{myblue}{Data Deletion:}
			\begin{itemize}
				\item Delete when no longer necessary
			\end{itemize}
			\item \textcolor{myblue}{Data Security:}
			\begin{itemize}
				\item Encryption, anonymization, differential privacy
			\end{itemize}
		\end{enumerate}
		
		\textcolor{myred}{\textbf{EXAM TRAP:}} Multiple principles guide \textcolor{myblue}{ethical privacy practices}!
	\end{frame}
	
	\begin{frame}{T94: Ethical Frameworks Summary}
		\fontsize{6.5}{7.5}\selectfont
		\textbf{TOPIC: Summary of Ethical Frameworks}
		
		\vspace{0.1cm}
		\textbf{CONSEQUENTIALISM:}
		\begin{itemize}
			\item Focus: Outcomes and consequences
			\item Key Question: "What maximizes good outcomes?"
			\item Strength: Single quantifiable metric
			\item Weakness: Can justify rights violations
		\end{itemize}
		
		\vspace{0.1cm}
		\textbf{DEONTOLOGY:}
		\begin{itemize}
			\item Focus: Duties and rules
			\item Key Question: "What are our moral obligations?"
			\item Strength: Respects rights and dignity
			\item Weakness: Can be rigid
		\end{itemize}
		
		\vspace{0.1cm}
		\textbf{VIRTUE ETHICS:}
		\begin{itemize}
			\item Focus: Character and virtues
			\item Key Question: "What kind of person should I be?"
			\item Strength: Holistic and aspirational
			\item Weakness: Lack of clear guidance
		\end{itemize}
		
		\textcolor{myred}{\textbf{EXAM TRAP:}} Each framework offers \textcolor{myblue}{valuable insights}; best decisions integrate all three!
	\end{frame}
	
	\begin{frame}{T95: AI Ethics Principles Summary}
		\fontsize{6.5}{7.5}\selectfont
		\textbf{TOPIC: Summary of AI Ethics Principles}
		
		\vspace{0.1cm}
		\textbf{FIVE CORE PRINCIPLES:}
		\begin{enumerate}
			\item \textcolor{myblue}{Nonmaleficence:}
			\begin{itemize}
				\item "Do no harm"
				\item Avoid unjustified risks
				\item Proportional risk-benefit
			\end{itemize}
			\item \textcolor{myblue}{Beneficence:}
			\begin{itemize}
				\item "Do good"
				\item Active help
				\item Advance welfare
			\end{itemize}
			\item \textcolor{myblue}{Justice:}
			\begin{itemize}
				\item Fairness and equality
				\item Impartiality
				\item Non-discrimination
			\end{itemize}
			\item \textcolor{myblue}{Autonomy:}
			\begin{itemize}
				\item Informed decisions
				\item Non-coercion
				\item Respect for persons
			\end{itemize}
			\item \textcolor{myblue}{Explainability:}
			\begin{itemize}
				\item Transparency
				\item Understandable decisions
				\item Accountability
			\end{itemize}
		\end{enumerate}
		
		\textcolor{myred}{\textbf{EXAM TRAP:}} Five principles provide \textcolor{myblue}{comprehensive ethical framework} for AI!
	\end{frame}
	
	\begin{frame}{T96: Governance Challenges Summary}
		\fontsize{6.5}{7.5}\selectfont
		\textbf{TOPIC: Summary of AI Governance Challenges}
		
		\vspace{0.1cm}
		\textbf{KEY CHALLENGES:}
		\begin{enumerate}
			\item \textcolor{myblue}{Power Asymmetries:}
			\begin{itemize}
				\item Tech companies more powerful than governments
			\end{itemize}
			\item \textcolor{myblue}{Lack of Transparency:}
			\begin{itemize}
				\item Private development, limited information
			\end{itemize}
			\item \textcolor{myblue}{Lack of Interpretability:}
			\begin{itemize}
				\item Black-box models, proxy variables
			\end{itemize}
			\item \textcolor{myblue}{Lack of Ethics Structures:}
			\begin{itemize}
				\item No required education, licensing, committees
			\end{itemize}
			\item \textcolor{myblue}{Lack of Regulation:}
			\begin{itemize}
				\item Few national laws, no specialized agencies
			\end{itemize}
			\item \textcolor{myblue}{Unpredictability:}
			\begin{itemize}
				\item Emergent behavior, unexpected capabilities
			\end{itemize}
			\item \textcolor{myblue}{Truth-Tracking:}
			\begin{itemize}
				\item LLMs give plausible, not truthful, responses
			\end{itemize}
		\end{enumerate}
		
		\textcolor{myred}{\textbf{EXAM TRAP:}} AI governance faces \textcolor{myblue}{multiple interconnected challenges}!
	\end{frame}
	
	\begin{frame}{T97: Regulatory Landscape Summary}
		\fontsize{6.5}{7.5}\selectfont
		\textbf{TOPIC: Summary of Global Regulatory Landscape}
		
		\vspace{0.1cm}
		\textbf{EUROPEAN UNION:}
		\begin{itemize}
			\item GDPR: Comprehensive privacy regulation
			\item DMA: Fair competition in digital markets
			\item DSA: Illegal content, transparency, disinformation
			\item AI Act: Risk-based AI regulation
			\item European AI Office: Enforcement body
		\end{itemize}
		
		\vspace{0.1cm}
		\textbf{UNITED STATES:}
		\begin{itemize}
			\item No federal privacy law
			\item State-level privacy laws (CCPA, CDPA, CPA, BIPA)
			\item Executive Order on AI
			\item NIST AI RMF (voluntary)
			\item AI Safety and Security Board
		\end{itemize}
		
		\vspace{0.1cm}
		\textbf{CHINA:}
		\begin{itemize}
			\item Generative AI Measures
			\item PIPL (privacy)
			\item Algorithmic Recommendation Rules
		\end{itemize}
		
		\textcolor{myred}{\textbf{EXAM TRAP:}} Different regions have \textcolor{myblue}{different regulatory approaches} - EU most comprehensive!
	\end{frame}
	
	\begin{frame}{T98: Practical Ethics Summary}
		\fontsize{6.5}{7.5}\selectfont
		\textbf{TOPIC: Summary of Practical Ethics}
		
		\vspace{0.1cm}
		\textbf{KEY ASPECTS:}
		\begin{itemize}
			\item Applied focus on actual ethical dilemmas
			\item Multidisciplinary approach
			\item Context-sensitive emphasis on situational nuances
			\item Pluralistic blend of different frameworks
		\end{itemize}
		
		\vspace{0.1cm}
		\textbf{BUSINESS BENEFITS:}
		\begin{itemize}
			\item Building trust and reputation
			\item Avoiding scandals
			\item Attracting and retaining talent
			\item Strengthening culture
			\item Supporting risk management
			\item Guiding innovation
			\item Promoting long-term thinking
		\end{itemize}
		
		\vspace{0.1cm}
		\textbf{KEY INSIGHT:}
		\begin{itemize}
			\item Practical ethics makes good business sense
			\item Benefits overlap and reinforce each other
			\item Ethics frameworks underlay responsible AI
		\end{itemize}
		
		\textcolor{myred}{\textbf{EXAM TRAP:}} Practical ethics = \textcolor{myblue}{applying theory to real-world dilemmas} with multiple benefits!
	\end{frame}
	
	\begin{frame}{T99: Responsible AI Implementation}
		\fontsize{6.5}{7.5}\selectfont
		\textbf{TOPIC: Implementing Responsible AI}
		
		\vspace{0.1cm}
		\textbf{KEY STEPS:}
		\begin{enumerate}
			\item \textcolor{myblue}{Establish Ethical Foundations:}
			\begin{itemize}
				\item Adopt ethical frameworks
				\item Define core principles
				\item Create ethics committees
			\end{itemize}
			\item \textcolor{myblue}{Implement Bias Mitigation:}
			\begin{itemize}
				\item Identify sources
				\item Use appropriate tools
				\item Regular auditing
			\end{itemize}
			\item \textcolor{myblue}{Protect Privacy:}
			\begin{itemize}
				\item Data minimization
				\item Strong security
				\item User control
			\end{itemize}
			\item \textcolor{myblue}{Ensure Explainability:}
			\begin{itemize}
				\item Transparent decisions
				\item Understandable explanations
				\item Accountability
			\end{itemize}
			\item \textcolor{myblue}{Establish Governance:}
			\begin{itemize}
				\item Oversight structures
				\item Continuous monitoring
				\item Stakeholder engagement
			\end{itemize}
		\end{enumerate}
		
		\textcolor{myred}{\textbf{EXAM TRAP:}} Responsible AI requires \textcolor{myblue}{comprehensive implementation} across all areas!
	\end{frame}
	
	\begin{frame}{T100: Responsible AI Summary}
		\fontsize{6.5}{7.5}\selectfont
		\textbf{TOPIC: Comprehensive Summary - Responsible \& Ethical AI}
		
		\vspace{0.1cm}
		\textbf{ETHICAL FOUNDATIONS:}
		\begin{itemize}
			\item Consequentialism (outcomes and utility)
			\item Deontology (duties and rules)
			\item Virtue Ethics (character and virtues)
		\end{itemize}
		
		\vspace{0.1cm}
		\textbf{CORE PRINCIPLES:}
		\begin{itemize}
			\item Nonmaleficence, beneficence, justice, autonomy, explainability
		\end{itemize}
		
		\vspace{0.1cm}
		\textbf{KEY PRACTICES:}
		\begin{itemize}
			\item Bias identification and mitigation
			\item Privacy protection and data minimization
			\item Transparency and explainability
			\item Human oversight and accountability
		\end{itemize}
		
		\vspace{0.1cm}
		\textbf{GOVERNANCE AND REGULATION:}
		\begin{itemize}
			\item Ethical committees and structures
			\item Algorithmic auditing
			\item Compliance with regulations (GDPR, EU AI Act)
			\item Continuous monitoring and improvement
		\end{itemize}
		
		\textcolor{myred}{\textbf{EXAM SUCCESS:}} Understand \textcolor{myblue}{ethical frameworks, principles, governance, and regulation} for responsible AI!
	\end{frame}
	
\end{document}